\documentclass[11pt,a4paper]{article}

% ── Codificación, fuentes y lenguaje ───────────────────────────────────────────
\usepackage[utf8]{inputenc}
\usepackage[T1]{fontenc}
\usepackage[default,scale=0.95]{sourcesanspro}
\renewcommand{\familydefault}{\sfdefault}
\usepackage[spanish,es-noshorthands,es-tabla]{babel}

% ── Geometría y espaciado ─────────────────────────────────────────────────────
\usepackage[top=2.2cm, bottom=2.5cm, left=2.2cm, right=2.2cm, headheight=15pt]{geometry}
\usepackage{setspace}
\setstretch{1.18}
\usepackage{parskip}

% ── Matemáticas y unidades ────────────────────────────────────────────────────
\usepackage{amsmath}
\usepackage{amssymb}
\usepackage{siunitx}

% ── Circuitos ─────────────────────────────────────────────────────────────────
\usepackage[european, straightvoltages]{circuitikz}

% ── Tablas ────────────────────────────────────────────────────────────────────
\usepackage{booktabs}
\usepackage{tabularx}
\usepackage{array}
\usepackage{longtable}
\usepackage{colortbl}
\usepackage{enumitem}

% ── Colores, cajas e iconos ───────────────────────────────────────────────────
\usepackage[dvipsnames]{xcolor}
\usepackage{tcolorbox}
\tcbuselibrary{skins, breakable}
\usepackage{fontawesome5}
\usepackage{graphicx}

% ── Listados de código (dark-mode infográfico) ────────────────────────────────
%   Estilo base: fondo oscuro con números de línea. Los bloques lstlisting se
%   envuelven automáticamente en una caja tcolorbox redondeada.
\usepackage{listings}

% ── Secciones con barra de color (infografía) ────────────────────────────────
\usepackage{titlesec}
\titleformat{\section}
  {\normalfont\LARGE\bfseries\color{azulNoche}}
  {\colorbox{cyanNeon}{\color{fondoOscuro}\thesection}}{0.7em}{}
  [\vspace{2pt}\color{cyanNeon}\rule{\linewidth}{1.8pt}]
\titlespacing*{\section}{0pt}{24pt}{10pt}
\titleformat{\subsection}
  {\normalfont\large\bfseries\color{azulNoche}}
  {\textcolor{cyanNeon}{\thesubsection}}{0.7em}{}
  [\vspace{1pt}\color{grisLinea}\rule{\linewidth}{0.6pt}]
\titlespacing*{\subsection}{0pt}{14pt}{6pt}
\titleformat{\subsubsection}
  {\normalfont\normalsize\bfseries\color{azulNoche}}
  {\textcolor{cyanNeon}{\thesubsubsection}}{0.7em}{}
\titlespacing*{\subsubsection}{0pt}{10pt}{4pt}

% ── Cabeceras y pies de página ────────────────────────────────────────────────
\usepackage{fancyhdr}
\usepackage{lastpage}
\pagestyle{fancy}
\fancyhf{}
\renewcommand{\headrulewidth}{0pt}
\renewcommand{\footrulewidth}{0pt}
\fancyfoot[C]{%
  \begin{minipage}{\linewidth}
    \vspace{2pt}
    \color{cyanNeon}\rule{\linewidth}{1.2pt}\\[2pt]
    \centering\color{grisTexto}\footnotesize
    Página \thepage\ de \pageref{LastPage}
  \end{minipage}%
}

% ── Hiperenlaces ──────────────────────────────────────────────────────────────
\usepackage[colorlinks=true, linkcolor=azulNoche, urlcolor=cyanNeon]{hyperref}
\sloppy
\emergencystretch 3em

% ── Paleta de colores — tecnológico dark-mode ──────────────────────────────────
\definecolor{fondoOscuro}{HTML}{0D1B2A}     % Dark navy — fondo banda título
\definecolor{verdeTurquesa}{HTML}{004D40}    % Verde turquesa — bandas de marca
\definecolor{azulNoche}{HTML}{1B3A6B}       % Midnight blue — secciones, marcos
\definecolor{azulMedio}{HTML}{2A5298}       % Azul medio — variante de acento
\definecolor{cyanNeon}{HTML}{00C8E0}        % Cyan eléctrico — acentos primarios
\definecolor{verdeSignal}{HTML}{00B887}     % Verde señal — fortalezas, éxito
\definecolor{naranjaVivo}{HTML}{FF7043}     % Naranja vivo — mejoras, advertencias
\definecolor{rojoAlerta}{HTML}{E53935}      % Rojo alerta — errores
\definecolor{grisPapel}{HTML}{F4F6FA}       % Fondo tarjetas claras
\definecolor{grisLinea}{HTML}{C8D0E0}       % Bordes sutiles
\definecolor{grisTexto}{HTML}{6B7A99}       % Texto secundario
\definecolor{amarilloNota}{HTML}{FFB300}    % Notas / advertencias doradas

% Colores específicos para bloques de código (dark-mode)
\definecolor{codigoFondo}{HTML}{1E2D3D}      % Fondo del bloque de código
\definecolor{codigoComentario}{HTML}{7A8FA6} % Comentarios
\definecolor{codigoCadena}{HTML}{FF9D5C}     % Cadenas / strings
\definecolor{codigoKeyword}{HTML}{00C8E0}    % Palabras clave
\definecolor{codigoNumero}{HTML}{00E5A0}     % Números

% ── Banda de título: portada infográfica ──────────────────────────────────────
%   Diseño en 2 capas: banda de color (por defecto fondo oscuro) + franja de
%   acento cyan abajo. Uso: \bandaTitulo[color]{título}{subtítulo}.
%   El icono se puede cambiar con \renewcommand{\iconoBanda}{...} (FontAwesome).
\newcommand{\iconoBanda}{microchip}
\newcommand{\bandaTitulo}[3][fondoOscuro]{%
  \begin{tcolorbox}[
    enhanced,
    colback=#1, colframe=#1,
    boxrule=0pt, arc=8pt, outer arc=8pt,
    width=\linewidth,
    top=18pt, bottom=6pt, left=14pt, right=14pt,
    borderline south={3.5pt}{0pt}{cyanNeon},
  ]
    \begin{minipage}[c]{0.07\linewidth}
      \centering
      \textcolor{cyanNeon}{\fontsize{30}{30}\selectfont\faIcon{\iconoBanda}}
    \end{minipage}%
    \hspace{8pt}%
    \begin{minipage}[c]{0.88\linewidth}
      {\fontsize{20}{24}\selectfont\bfseries\color{white}#2}\par\vspace{5pt}%
      \textcolor{cyanNeon!80!white}{\small\faIcon{angle-right}~#3}%
    \end{minipage}
    \vspace{4pt}
  \end{tcolorbox}%
  \vspace{8pt}%
}

% ── Tarjetas de datos (infografía) ────────────────────────────────────────────
\tcbset{
  tarjetaDato/.style={
    enhanced, breakable,
    arc=6pt, outer arc=6pt,
    colback=grisPapel, colframe=azulNoche,
    boxrule=1pt,
    fonttitle=\bfseries\small, coltitle=white,
    attach boxed title to top left={yshift=-3mm, xshift=6mm},
    boxed title style={
      colback=azulNoche, colframe=azulNoche,
      arc=4pt, boxrule=0pt,
    },
    drop fuzzy shadow=azulNoche!30!white,
    top=8pt, bottom=8pt, left=10pt, right=10pt,
  },
  tarjetaIcono/.style={
    enhanced, breakable,
    arc=6pt, outer arc=6pt,
    colback=white, colframe=grisLinea, boxrule=0.8pt,
    fonttitle=\bfseries\small, coltitle=azulNoche,
    borderline west={3pt}{0pt}{cyanNeon},
    drop fuzzy shadow=azulNoche!25!white,
    top=6pt, bottom=6pt, left=10pt, right=10pt,
  },
  cajaTitulo/.style={
    enhanced, breakable,
    arc=6pt, outer arc=6pt,
    colback=azulNoche!10!grisPapel, colframe=azulNoche,
    boxrule=1pt,
    fonttitle=\bfseries, coltitle=white,
    attach boxed title to top left={yshift=-3mm, xshift=6mm},
    boxed title style={
      colback=azulNoche, colframe=azulNoche,
      arc=4pt, boxrule=0pt,
    },
    drop fuzzy shadow=azulNoche!25!white,
    top=8pt, bottom=8pt, left=10pt, right=10pt,
  },
  cajaContenido/.style={
    enhanced, breakable,
    arc=5pt, outer arc=5pt,
    colback=grisPapel, colframe=grisLinea,
    boxrule=0.8pt,
    fonttitle=\bfseries\small, coltitle=azulNoche,
    top=6pt, bottom=6pt, left=10pt, right=10pt,
  },
  cajaFortaleza/.style={
    enhanced, breakable,
    arc=6pt, outer arc=6pt,
    colback=verdeSignal!8!white, colframe=verdeSignal,
    boxrule=1pt,
    borderline west={4pt}{0pt}{verdeSignal},
    fonttitle=\bfseries\small, coltitle=verdeSignal!70!black,
    drop fuzzy shadow=verdeSignal!20!white,
    top=6pt, bottom=6pt, left=10pt, right=10pt,
  },
  cajaMejora/.style={
    enhanced, breakable,
    arc=6pt, outer arc=6pt,
    colback=naranjaVivo!8!white, colframe=naranjaVivo,
    boxrule=1pt,
    borderline west={4pt}{0pt}{naranjaVivo},
    fonttitle=\bfseries\small, coltitle=naranjaVivo!80!black,
    drop fuzzy shadow=naranjaVivo!20!white,
    top=6pt, bottom=6pt, left=10pt, right=10pt,
  },
  cajaRecomendacion/.style={
    enhanced, breakable,
    arc=6pt, outer arc=6pt,
    colback=cyanNeon!6!white, colframe=cyanNeon!60!azulNoche,
    boxrule=1pt,
    borderline west={4pt}{0pt}{cyanNeon},
    fonttitle=\bfseries\small, coltitle=azulNoche,
    drop fuzzy shadow=azulNoche!20!white,
    top=6pt, bottom=6pt, left=10pt, right=10pt,
  },
}

% ── Ícono + texto (encabezados de sección y elementos) ────────────────────────
\newcommand{\iconotexto}[2]{\textcolor{cyanNeon}{\faIcon{#1}}~#2}

% ── Listados de código: estilo dark + auto-envoltura ─────────────────────────
\lstdefinestyle{estiloCodigo}{
    backgroundcolor=\color{codigoFondo},
    basicstyle=\ttfamily\small\color{white},
    commentstyle=\color{codigoComentario}\itshape,
    stringstyle=\color{codigoCadena},
    keywordstyle=\color{codigoKeyword}\bfseries,
    numberstyle=\tiny\color{codigoComentario},
    numbers=left,
    numbersep=10pt,
    showstringspaces=false,
    breaklines=true,
    breakatwhitespace=false,
    tabsize=4,
    frame=none,
    captionpos=b,
    aboveskip=6pt,
    belowskip=4pt,
    xleftmargin=14pt,
    framexleftmargin=14pt,
    literate=
      {á}{{\'a}}1 {é}{{\'e}}1 {í}{{\'i}}1 {ó}{{\'o}}1 {ú}{{\'u}}1
      {Á}{{\'A}}1 {É}{{\'E}}1 {Í}{{\'I}}1 {Ó}{{\'O}}1 {Ú}{{\'U}}1
      {ñ}{{\~n}}1 {Ñ}{{\~N}}1 {ü}{{\"u}}1 {Ü}{{\"U}}1
      {¿}{{?`}}1 {¡}{{!`}}1
}
\lstdefinestyle{estiloPython}{
    style=estiloCodigo,
    language=Python,
}
\lstset{style=estiloCodigo}
\BeforeBeginEnvironment{lstlisting}{%
  \begin{tcolorbox}[
    enhanced, breakable,
    arc=6pt, outer arc=6pt,
    colback=codigoFondo, colframe=azulNoche,
    boxrule=1pt,
    drop fuzzy shadow=azulNoche!25!white,
    top=2pt, bottom=2pt, left=0pt, right=0pt,
  ]%
}
\AfterEndEnvironment{lstlisting}{\end{tcolorbox}}

% ── Cajas de contenido temáticas ──────────────────────────────────────────────
\newtcolorbox{cajaRecuerda}{
    enhanced, breakable,
    arc=6pt, outer arc=6pt,
    colback=azulNoche!10!grisPapel,
    colframe=azulNoche, boxrule=1pt,
    borderline west={4pt}{0pt}{cyanNeon},
    fonttitle=\bfseries\small\color{white},
    title={\faIcon{info-circle}~Recuerda},
    attach boxed title to top left={yshift=-3mm, xshift=6mm},
    boxed title style={colback=azulNoche, arc=4pt, boxrule=0pt},
    drop fuzzy shadow=azulNoche!25!white,
    left=10pt, right=10pt, top=8pt, bottom=8pt,
}

\newtcolorbox{cajaConcepto}{
    enhanced, breakable,
    arc=6pt, outer arc=6pt,
    colback=verdeSignal!8!white,
    colframe=verdeSignal, boxrule=1pt,
    borderline west={4pt}{0pt}{verdeSignal},
    fonttitle=\bfseries\small\color{white},
    title={\faIcon{lightbulb}~Concepto clave},
    attach boxed title to top left={yshift=-3mm, xshift=6mm},
    boxed title style={colback=verdeSignal!80!black, arc=4pt, boxrule=0pt},
    drop fuzzy shadow=verdeSignal!20!white,
    left=10pt, right=10pt, top=8pt, bottom=8pt,
}

\newtcolorbox{cajaEjemplo}{
    enhanced, breakable,
    arc=6pt, outer arc=6pt,
    colback=cyanNeon!5!white,
    colframe=cyanNeon!70!azulNoche, boxrule=1pt,
    borderline west={4pt}{0pt}{cyanNeon},
    fonttitle=\bfseries\small\color{fondoOscuro},
    title={\faIcon{code}~Ejemplo},
    attach boxed title to top left={yshift=-3mm, xshift=6mm},
    boxed title style={colback=cyanNeon, arc=4pt, boxrule=0pt},
    drop fuzzy shadow=azulNoche!20!white,
    left=10pt, right=10pt, top=8pt, bottom=8pt,
}

\newtcolorbox{cajaEjercicio}{
    enhanced, breakable,
    arc=6pt, outer arc=6pt,
    colback=azulMedio!7!white,
    colframe=azulMedio, boxrule=1pt,
    borderline west={4pt}{0pt}{azulMedio},
    fonttitle=\bfseries\small\color{white},
    title={\faIcon{pencil-alt}~Ejercicio},
    attach boxed title to top left={yshift=-3mm, xshift=6mm},
    boxed title style={colback=azulMedio, arc=4pt, boxrule=0pt},
    drop fuzzy shadow=azulNoche!20!white,
    left=10pt, right=10pt, top=8pt, bottom=8pt,
}

\newtcolorbox{cajaReto}{
    enhanced, breakable,
    arc=6pt, outer arc=6pt,
    colback=naranjaVivo!6!white,
    colframe=naranjaVivo, boxrule=1pt,
    borderline west={4pt}{0pt}{naranjaVivo},
    fonttitle=\bfseries\small\color{white},
    title={\faIcon{fire}~Reto},
    attach boxed title to top left={yshift=-3mm, xshift=6mm},
    boxed title style={colback=naranjaVivo, arc=4pt, boxrule=0pt},
    drop fuzzy shadow=naranjaVivo!20!white,
    left=10pt, right=10pt, top=8pt, bottom=8pt,
}

% ── Indicadores de dificultad ─────────────────────────────────────────────────
\newcommand{\facil}{\textcolor{verdeSignal}{\faIcon{circle}\,\textbf{Fácil}}}
\newcommand{\intermedio}{\textcolor{naranjaVivo}{\faIcon{circle}\,\textbf{Intermedio}}}
\newcommand{\dificil}{\textcolor{rojoAlerta}{\faIcon{circle}\,\textbf{Difícil}}}

% ─────────────────────────────────────────────────────────────────────────────

% ── Cabecera específica del reporte del skill ──
\renewcommand{\iconoBanda}{book}
\fancyhead[L]{\color{azulNoche}\small\bfseries \faIcon{book}~\textcolor{cyanNeon}{dsh}}
\fancyhead[R]{\color{grisTexto}\smallAgent harness: Cordis plugin kernel, 4 runtime modes, trazabilidad append-only, instalación y operación}

\begin{document}
\bandaTitulo[fondoOscuro]{Skill: dsh}{"Referencia del Agent Harness DeepSeek (dsh) — kernel de plugins Cordis, 4 modos runtime, trazabilidad append-only, instalación y operación. Úsalo cuando trabajes con el harness dsh, perfiles, bundles, plugins Cordis, ev}

\section*{dsh — Agent Harness: Cordis Plugin Kernel}


\subsection*{Cuándo usar este skill}


Usa este skill cuando necesites:


\begin{itemize}[leftmargin=*,topsep=2pt,itemsep=1pt]
  \item Entender la arquitectura del harness dsh (kernel Cordis, seams, servicios core).
  \item Trabajar con perfiles (\texttt{dsh --profile}), bundles, parches y overlays.
  \item Gestionar plugins (\texttt{dsh plugin}), subagentes (Codex, Claude Code, ACP, SDK).
  \item Implementar servicios remotos con Typert (\texttt{@Remote}, \texttt{@RemoteScope}).
  \item Consumir eventos del ciclo de vida del agente (\texttt{agent/\textit{}, \texttt{session/}}, \texttt{tools/*}).
  \item Configurar plugins en \texttt{cordis.yml} con sus interfaces \texttt{Config}.
  \item Operar los 4 modos runtime: web, headless, sdk, acp.
\end{itemize}


\subsection*{Estructura del proyecto}


El harness se organiza en paquetes npm bajo el scope \texttt{@deepseek-ai/}:


\begin{itemize}[leftmargin=*,topsep=2pt,itemsep=1pt]
  \item \textbf{Core}: \texttt{dsh-agent}, \texttt{dsh-agent-loop}, \texttt{dsh-session}, \texttt{dsh-tools}, \texttt{dsh-system-prompt} — servicios dueños del estado.
  \item \textbf{Seams}: \texttt{dsh-llm}, \texttt{dsh-fs}, \texttt{dsh-storage}, \texttt{dsh-subagents}, \texttt{dsh-compaction}, \texttt{dsh-sandbox} — interfaces abstractas con múltiples backends.
  \item \textbf{Bundles}: \texttt{dsh-base}, \texttt{dsh-web-app}, \texttt{dsh-headless}, \texttt{dsh-sdk-app}, \texttt{dsh-sdk-minimal}, \texttt{dsh-acp-app} — composiciones de aplicación.
  \item \textbf{Protocolo}: \texttt{dsh-typert-protocol} — RPC tipado entre Host y Cliente.
  \item \textbf{Frontend}: \texttt{dsh-client-*}, \texttt{dsh-host-webserver}, \texttt{dsh-web-app}.
\end{itemize}


\subsection*{API principal}


\subsubsection*{Contexto Cordis (\texttt{ctx})}


\begin{lstlisting}
// Ciclo de vida
ctx.extend(meta?: {}): this
ctx.isolate(name: string, label?: symbol): this
ctx.intercept(name: string, config: any): this

// Servicios
ctx.get(name: string, strict?: boolean): any
ctx.set(name: string, value: any): void
ctx.provide(name: string, value?: any): () => void
ctx.accessor(name: string, options: { get?, set? }): void
ctx.mixin(name: string, mixins: string[] | Dict<string>): void

// Plugins
ctx.plugin(plugin: Plugin, ...args: any[]): Fiber & PromiseLike<Fiber>
ctx.inject(deps: Inject, callback: Plugin.Function<void>): Fiber & PromiseLike<Fiber>

// Efectos
ctx.effect(execute: () => SyncEffect, label?: string): Disposable<Promise<void>>
ctx.effect(execute: () => Effect, label?: string): AsyncDisposable<Promise<void>>
\end{lstlisting}


\subsubsection*{Eventos}


\begin{lstlisting}
ctx.parallel(name: string, ...args: any[]): Promise<void>  // concurrente
ctx.emit(name: string, ...args: any[]): void                // síncrono
ctx.serial(name: string, ...args: any[]): Promise<any>      // secuencial async
ctx.bail(name: string, ...args: any[]): any                 // secuencial sync
ctx.waterfall(name: string, ...args: any[]): any            // middleware encadenado
ctx.on(name: string, listener: Function, options?: boolean | EventOptions): () => boolean
ctx.once(name: string, listener: Function, options?: boolean | EventOptions): () => boolean
\end{lstlisting}


\subsubsection*{Typert Remote}


\begin{lstlisting}
class TypertRemoteService {
  constructor(ctx: Context, namespace: string)
}

function bindTypertRemote(service: any, serviceKey: string): void

// Decoradores
@Remote(methodName: string)          // endpoint RPC unario
@RemoteScope(key: string, methodName: string)  // RPC con ámbito (ej. 'agent')

// Cliente
ctx.remote.<namespace>.<method>(...)
agentCtx.remote.<namespace>.<method>(...)
\end{lstlisting}


\subsubsection*{Service (clase base)}


\begin{lstlisting}
class Service {
  public name!: string
  constructor(ctx: Context, name: string)
  static readonly init: unique symbol
  static readonly check: unique symbol
  static readonly config: unique symbol
  static readonly invoke: unique symbol
  static readonly extend: unique symbol
  static readonly tracker: unique symbol
  static readonly resolveConfig: unique symbol
}
\end{lstlisting}


\subsubsection*{Fiber}


\begin{lstlisting}
fiber.uid: number | null
fiber.ctx: Context
fiber.config: any
fiber.state: FiberState  // PENDING → LOADING → ACTIVE → UNLOADING → DISPOSED
fiber.dispose(): Promise<void>
fiber.name: string
fiber.store: Dict<Impl> | undefined
fiber.inertia: Promise<void> | undefined
fiber.assertActive(): void
fiber.getEffects(): EffectMeta[]
fiber.await(): Promise<Fiber>
fiber.restart(): Promise<void>
fiber.update(config: any, noSave?: boolean): any
\end{lstlisting}


\subsubsection*{Herramientas}


\begin{lstlisting}
ctx.tools.register(toolDefinition): Disposer
ctx.tools.restrict(scopedRegistration): void
ctx.tools.guard(exec => boolean): void
ctx.tools.execute(params: {
  callId: ToolCallId
  name: string
  arguments: Record<string, unknown>
  signal: AbortSignal
}): Promise<{ content: ContentBlock[] }>
\end{lstlisting}


\subsection*{Patrones de uso}


\subsubsection*{1. Registrar un servicio}


\begin{lstlisting}
declare module '@deepseek-ai/cordis' {
  interface Context { greeter: GreeterService }
}

export class GreeterService extends Service {
  constructor(ctx: Context) { super(ctx, 'greeter') }
  greet(who: string) { return `Hello, ${who}!` }
}

export function apply(ctx: Context) { ctx.plugin(GreeterService) }
\end{lstlisting}


\subsubsection*{2. Consumir con dependencia dura}


\begin{lstlisting}
export const inject = ['greeter']
export function apply(ctx: Context) {
  console.log(ctx.greeter.greet('world'))
}
\end{lstlisting}


\subsubsection*{3. Waterfall para decisiones}


\begin{lstlisting}
ctx.on('some-decision-event', function (payload, next) {
  if (payload.blocked) return { denied: true }  // cortocircuito
  return next()
})
\end{lstlisting}


\subsubsection*{4. Servicio remoto Typert}


\begin{lstlisting}
export class GoalService extends TypertRemoteService {
  constructor(ctx: Context) { super(ctx, 'goals') }

  @Remote('create')
  createForClient(agent: Agent, request: { objective: string }, signal: AbortSignal) {
    signal.throwIfAborted()
    return this.create(agent, request)
  }

  @RemoteScope('agent', 'current')
  currentForClient() { return { accepted: true } }
}
\end{lstlisting}


\subsubsection*{5. Consumo remoto con manejo de errores}


\begin{lstlisting}
export const inject = ['remote', 'remote.notes']

const result = await ctx.remote.notes.list()
if (!result.ok) {
  if (result.error.code === 'note/not-found') return []
  throw result.error
}
return result.value
\end{lstlisting}


\subsection*{Configuración}


\subsubsection*{Perfiles}


\begin{lstlisting}
dsh --profile web          # base + web-app, parches en vivo
dsh --profile headless     # base + headless, solo startup
dsh --profile sdk          # base + sdk-app, solo startup
dsh --profile sdk-minimal  # bundle independiente
dsh --profile acp          # base + acp-app, solo startup
\end{lstlisting}


\subsubsection*{Capas de configuración (orden de aplicación)}


\begin{enumerate}[leftmargin=*,topsep=2pt,itemsep=1pt]
  \item Parches de bundles en \texttt{dsh.profile.bundles}
  \item \texttt{cordis.patch.yml} del perfil
  \item \texttt{$DSH\_HOME/cordis.patch.yml}
  \item Overlays \texttt{--patch <path>} en orden
\end{enumerate}


\subsubsection*{Gestión de plugins}


\begin{lstlisting}
dsh plugin --profile <name> add <package>
dsh plugin --profile <name> remove <package>
dsh plugin --profile <name> add github:deepseek-harness/turtle-ui
\end{lstlisting}


\subsubsection*{cordis.yml}


\begin{lstlisting}
plugins:
  '@deepseek-ai/dsh-agent-loop':
    config:
      maxParallelToolCalls: 4
      agents:
        - id: main-agent
          provider: deepseek
          model: deepseek-chat
          cwd: /home/user/project
\end{lstlisting}


\subsection*{Ejemplos de código reutilizable}


Ver \texttt{references/ejemplos.md} para fragmentos completos.


\subsection*{Glosario}


Ver \texttt{references/glosario.md} para términos técnicos.


\subsection*{Prerrequisitos}


Ver \texttt{references/prerrequisitos.md} para dependencias e instalación.


\section*{api}

\section*{API principal del proyecto dsh}


\subsection*{Contexto Cordis (\texttt{ctx})}


Objeto central: servicios, eventos y ciclo de vida accesibles a través de \texttt{ctx}.


\subsubsection*{Métodos de ciclo de vida}


\begin{lstlisting}
// Crear hijo con metadatos adicionales (herencia prototípica)
ctx.extend(meta?: {}): this

// Crear hijo con ámbito de servicio independiente para 'name'
ctx.isolate(name: string, label?: symbol): this

// Añadir configuración de interceptación para plugins bajo este contexto
ctx.intercept(name: string, config: any): this
\end{lstlisting}


\subsubsection*{Propiedades}


\begin{lstlisting}
ctx.root        // Contexto raíz de la aplicación (@experimental)
ctx.baseUrl     // URL base para resolver especificadores relativos
ctx.events      // Servicio de bus de eventos
ctx.logger      // Servicio de logging; usar ctx.logger(name) para logger nombrado
ctx.reflect     // Capa de reflexión (ctx.get, ctx.provide, ...)
ctx.registry    // Registro de plugins (ctx.plugin, ctx.inject)
\end{lstlisting}


\subsubsection*{Miembros estáticos}


\begin{lstlisting}
Context.effect    // Symbol para árbol de diagnóstico EffectMeta en disposer
Context.filter    // Symbol para filtro de escucha del contexto
Context.isolate   // Symbol del mapa de aislamiento
Context.intercept // Symbol del mapa de interceptación
Context.is(value) // Type guard: true si value es un contexto Cordis
\end{lstlisting}


\subsubsection*{Almacén de servicios y mixins}


\begin{lstlisting}
ctx.get(name: string, strict?: boolean): any
ctx.set(name: string, value: any): void
ctx.provide(name: string, value?: any): () => void
ctx.accessor(name: string, options: { get?, set? }): void
ctx.mixin(name: string, mixins: string[] | Dict<string>): void
\end{lstlisting}


\subsection*{API de Eventos}


\subsubsection*{Métodos de despacho}


\begin{lstlisting}
ctx.parallel(name: string, ...args: any[]): Promise<void>  // concurrente
ctx.emit(name: string, ...args: any[]): void                // síncrono, ignora retorno
ctx.serial(name: string, ...args: any[]): Promise<any>      // secuencial async hasta bail
ctx.bail(name: string, ...args: any[]): any                 // secuencial sync hasta bail
ctx.waterfall(name: string, ...args: any[]): any            // middleware encadenado
\end{lstlisting}


\subsubsection*{Registro de listeners}


\begin{lstlisting}
ctx.on(name: string, listener: Function, options?: boolean | EventOptions): () => boolean
ctx.once(name: string, listener: Function, options?: boolean | EventOptions): () => boolean
\end{lstlisting}


\subsubsection*{Tipos}


\begin{lstlisting}
interface EventOptions {
  prepend?: boolean  // Añadir antes de listeners existentes
  global?: boolean   // Recibir evento ignorando filtros de contexto
}

type DispatchMode = 'emit' | 'parallel' | 'serial' | 'bail' | 'waterfall'
\end{lstlisting}


\subsection*{Typert Remote Service}


\subsubsection*{\texttt{TypertRemoteService}}


Clase base para servicios de negocio expuestos al cliente. El constructor enlaza explícitamente la clave de servicio Cordis y el namespace Remoto por defecto.


\begin{lstlisting}
constructor(ctx: Context, namespace: string)
\end{lstlisting}


\subsubsection*{\texttt{bindTypertRemote(service, serviceKey)}}


Alternativa para servicios con otra clase base. Declara un binding público inspeccionable.


\subsubsection*{Decoradores}


\begin{lstlisting}
@Remote(methodName: string)
// Decorador para exponer un método de instancia público y no estático como endpoint RPC unario.

@RemoteScope(key: string, methodName: string)
// Decorador para métodos que operan sobre un contexto resuelto por ámbito (ej. 'agent'),
// sin recibir el objeto explícitamente.
\end{lstlisting}


\subsubsection*{Interfaces de cliente}


\begin{lstlisting}
ctx.remote.<namespace>.<method>(...)       // llamadas RPC directas
agentCtx.remote.<namespace>.<method>(...)  // llamadas RPC con ámbito de agente
\end{lstlisting}


\subsubsection*{Manejo de errores remotos}


\begin{lstlisting}
// Verificar result.ok; si es false, ramificar por result.error.code (no instanceof)
// Lanzar result.error para propagar errores remotos
// Usar isRemoteFailure(error) para distinguir fallos remotos de defectos locales
// Un aborto resulta en error con código gateway/cancelled
\end{lstlisting}


\subsection*{Service (clase base)}


\begin{lstlisting}
class Service {
  public name!: string
  constructor(ctx: Context, name: string) // Registra la instancia en ctx.<name>

  static readonly init: unique symbol           // Método post-constructor (class plugins)
  static readonly check: unique symbol          // Predicado de disponibilidad para ctx.provide()
  static readonly config: unique symbol         // Phantom type para intercept-config
  static readonly invoke: unique symbol         // Cuerpo de llamada (ej. ctx.logger())
  static readonly extend: unique symbol         // Helper para derivar instancia extendida
  static readonly tracker: unique symbol        // Metadata de rastreo de contexto
  static readonly resolveConfig: unique symbol  // Helper de resolución de intercept-config
}
\end{lstlisting}


\subsection*{Fiber}


Instancia de plugin en runtime.


\begin{lstlisting}
fiber.uid: number | null          // ID único; 0 root, null si disposed
fiber.ctx: Context                // Contexto de ejecución del plugin
fiber.config: any                 // Configuración validada
fiber.state                       // Estado del ciclo de vida (emite 'internal/status')
fiber.dispose(): Promise<void>    // Descargar plugin y ejecutar cleanup
fiber.name: string                // Nombre display (hereda del ancestro, o 'root')
fiber.store: Dict<Impl> | undefined // Snapshot de servicios requeridos (mientras cargado)
fiber.inertia: Promise<void> | undefined // Transición load/unload en vuelo
\end{lstlisting}


\subsubsection*{Máquina de estados}


\begin{lstlisting}
PENDING → LOADING → ACTIVE → UNLOADING → DISPOSED
                 ↘ FAILED
\end{lstlisting}


\subsubsection*{Métodos adicionales}


\begin{lstlisting}
fiber.assertActive()              // Lanza CordisError si ya disposed
fiber.getEffects(): EffectMeta[]  // Árbol de efectos activos etiquetados
fiber.await(): Promise<Fiber>     // Espera fin de trabajo lifecycle; relanza errores de startup
fiber.restart(): Promise<void>    // Dispose + recarga inmediata con la config actual
fiber.update(config: any, noSave?: boolean): any // Valida, aplica y reinicia
\end{lstlisting}


\subsection*{Efectos}


\begin{lstlisting}
ctx.effect(execute: () => SyncEffect, label?: string): Disposable<Promise<void>>
ctx.effect(execute: () => Effect, label?: string): AsyncDisposable<Promise<void>>

// Tipos
type Effect<T = any> = SyncEffect<T> | AsyncEffect<T>
type Disposable<T = any> = () => T

// EffectMeta (diagnóstico)
interface EffectMeta { label: string, children: EffectMeta[] }
\end{lstlisting}


\subsection*{Manejo de errores}


\begin{lstlisting}
class CordisError extends Error {
  constructor(public code: CordisError.Code, message?: string)
}
// Códigos: { INACTIVE_EFFECT: 'cannot create effect on inactive context' }

class ValidationError extends TypeError {
  name = 'ValidationError'
  constructor(issues: readonly StandardSchemaV1.Issue[])
}
\end{lstlisting}


\subsection*{Plugin}


\begin{lstlisting}
ctx.plugin<P extends Plugin>(plugin: P, ...args: Spread<GetPluginConfig<P>>): Fiber & PromiseLike<Fiber>
ctx.inject(deps: Inject, callback: Plugin.Function<void>): Fiber & PromiseLike<Fiber>

// Formas de Plugin
type Plugin<T> = Plugin.Function<T> | Plugin.Constructor<T> | Plugin.Object<T>
// Plugin.Function<T>: (ctx: Context, config: T) => any  (con metadata Base)
// Plugin.Constructor<T>: new (ctx: Context, config: T) => any
// Plugin.Object<T>: { apply(ctx: Context, config: T): any }
// Base: { name?, Config?, inject?, provide?, intercept? }

// Inject
type Inject = (keyof M)[] | { [K in keyof M]?: M[K] }
\end{lstlisting}


\subsection*{Herramientas (harness)}


\begin{lstlisting}
ctx.tools.register(toolDefinition): Disposer
ctx.tools.restrict(scopedRegistration): void
ctx.tools.guard(exec => boolean): void
ctx.tools.execute(params: {
  callId: ToolCallId
  name: string
  arguments: Record<string, unknown>
  signal: AbortSignal
}): Promise<{ content: ContentBlock[] }>
\end{lstlisting}


\subsection*{Servicios core principales}


\begin{center}
\begin{tabular}{@{}lll@{}}
\toprule
Servicio & Dueño & Rol \\
\midrule
\texttt{ctx.sessions} & \texttt{session} & Almacén de sesiones en memoria, feed de eventos append-only \\
\texttt{ctx.tools} & \texttt{tools} & Registro de herramientas + pipeline de ejecución custodiado \\
\texttt{ctx.agents} & \texttt{agent} & Fábrica de agentes (create/resume), handles vivos \\
\texttt{ctx.systemPrompt} & \texttt{system-prompt} & Ensamblador de system prompt por paso \\
\texttt{ctx.agentLoop} & \texttt{agent-loop} & Driver concreto del bucle agente \\
\texttt{ctx.storageDomain} & \texttt{storage-domain} & Instalación tipada sobre backends de almacenamiento \\
\texttt{ctx.workspaceRegistry} & \texttt{workspace} & Registro de entidades Workspace con ID de marca \\
\texttt{ctx.invariants} & \texttt{invariants} & Checks invariantes con atribución de paquete \\
\texttt{ctx.planMode} & \texttt{plan-mode} & Estado de colaboración plan/mode plegado del log \\
\texttt{ctx.commands} & \texttt{commands} & Comandos directos del usuario sin pasar por modelo \\
\bottomrule
\end{tabular}
\end{center}



\subsection*{Seams reemplazables}


\begin{lstlisting}
// Seam ctx.llm — backends múltiples
// Dueño: llm
// Implementaciones: llm-deepseek, llm-pi-ai, llm-replay
// Consumidores: agent-loop, compaction-basic

// Seam ctx.fs — proveedor de filesystem
// Dueño: fs
// Implementaciones: fs-local, fs-sandbox, fs-e2b
// Consumidores: tool-fs (+ fs-observation-policy vía gate de eventos)
\end{lstlisting}


\subsection*{Eventos clave del ciclo de vida del agente}


\begin{center}
\begin{tabular}{@{}llll@{}}
\toprule
Evento & Modo & Dispatcher & Listeners principales \\
\midrule
\texttt{agent/created} & \texttt{emit} & \texttt{agent} & \texttt{agent-presets}, \texttt{file-reference-local}, \texttt{goal-round-driver}, \texttt{schedule}, \texttt{tool-agent-team}, \texttt{tool-subagent} \\
\texttt{agent/disposed} & \texttt{emit} & \texttt{agent} & \texttt{agent-loop}, \texttt{file-reference-local}, \texttt{goal-round-driver}, \texttt{session-controller}, \texttt{subagent}, \texttt{tool-agent-team}, \texttt{tool-subagent} \\
\texttt{agent/session-start} & \texttt{emit} & \texttt{agent-loop} & \texttt{agent-team}, \texttt{goal}, \texttt{goal-round-driver}, \texttt{hooks-claude-code}, \texttt{hooks-codex} \\
\texttt{agent/status} & \texttt{emit} & \texttt{agent-loop} & \texttt{agent}, \texttt{agent-team}, \texttt{compaction-basic}, \texttt{goal-round-driver}, \texttt{schedule}, \texttt{server}, \texttt{session-controller} \\
\texttt{agent/pre-step} & \texttt{waterfall} & \texttt{agent-loop} & \texttt{agent}, \texttt{agent-instructions}, \texttt{compaction-basic}, \texttt{goal-round-driver}, \texttt{hooks-claude-code}, \texttt{hooks-codex}, \texttt{plan-mode}, \texttt{repeat-tool-reminder}, \texttt{session-checkpoint-policy}, \texttt{session-reference}, \texttt{subagent-in-process-driver}, \texttt{time-context}, \texttt{tmux-context}, \texttt{tool-cordis}, \texttt{tool-skill}, \texttt{tool-subagent} \\
\texttt{agent/request} & \texttt{waterfall} & \texttt{agent-loop} & \texttt{agent}, \texttt{webhook} \\
\texttt{agent/request-error} & \texttt{waterfall} & \texttt{agent-loop} & \texttt{compaction-basic}, \texttt{llm-retry} \\
\texttt{agent/turn-stopping} & \texttt{serial} & \texttt{agent-loop} & \texttt{hooks-claude-code}, \texttt{hooks-codex} \\
\texttt{agent/error} & \texttt{emit} & \texttt{agent-loop} & \texttt{acp}, \texttt{goal-round-driver}, \texttt{session-controller}, \texttt{session-telemetry} \\
\texttt{agent/assistant-stream} & \texttt{emit} & \texttt{agent-loop} & \texttt{headless}, \texttt{session-controller} \\
\texttt{agent/inbox/inserted} & \texttt{emit} & \texttt{agent-loop} & \texttt{goal-round-driver} \\
\texttt{agent/inbox/claimed} & \texttt{emit} & \texttt{agent-loop} & \texttt{acp}, \texttt{goal-round-driver}, \texttt{subagent}, \texttt{tool-jobs} \\
\texttt{agent/inbox/discarded} & \texttt{emit} & \texttt{agent-loop} & \texttt{goal-round-driver}, \texttt{subagent} \\
\texttt{agent-loop/config-start-failed} & \texttt{emit} & \texttt{agent-loop} & - \\
\bottomrule
\end{tabular}
\end{center}



\subsection*{Eventos de sesión}


\begin{center}
\begin{tabular}{@{}llll@{}}
\toprule
Evento & Modo & Dispatcher & Listeners principales \\
\midrule
\texttt{session/created} & \texttt{emit} & \texttt{session} & \texttt{compaction}, \texttt{goal}, \texttt{hook-protocol}, \texttt{llm-retry}, \texttt{permission-presets}, \texttt{plan-mode}, \texttt{schedule}, \texttt{server}, \texttt{session}, \texttt{session-controller}, \texttt{session-log-deepseek}, \texttt{session-projection}, \texttt{session-projection-cache}, \texttt{session-telemetry}, \texttt{session-title}, \texttt{time-context}, \texttt{tool-todo}, \texttt{tool-workflow}, \texttt{tools}, \texttt{user-approval} \\
\texttt{session/disposed} & \texttt{emit} & \texttt{session} & \texttt{agent-loop}, \texttt{agent-team}, \texttt{file-upload}, \texttt{session-controller}, \texttt{session-persistence-jsonl}, \texttt{session-projection-cache}, \texttt{session-telemetry}, \texttt{session-title} \\
\texttt{session/event} & \texttt{emit} & \texttt{session} & \texttt{acp}, \texttt{agent-instructions}, \texttt{agent-loop}, \texttt{agent-presets}, \texttt{agent-team}, \texttt{compaction}, \texttt{compaction-basic}, \texttt{file-reference-local}, \texttt{file-upload}, \texttt{goal}, \texttt{goal-round-driver}, \texttt{hook-protocol}, \texttt{loader-smoke}, \texttt{server}, \texttt{session}, \texttt{session-controller}, \texttt{session-persistence-jsonl}, \texttt{session-projection}, \texttt{session-projection-cache}, \texttt{session-telemetry}, \texttt{session-telemetry-otel}, \texttt{session-title}, \texttt{token-meter}, \texttt{tool-todo}, \texttt{tool-workflow}, \texttt{tools}, \texttt{user-approval} \\
\texttt{session/flush} & \texttt{parallel} & \texttt{session} & \texttt{session-persistence-jsonl}, \texttt{session-telemetry} \\
\bottomrule
\end{tabular}
\end{center}



\subsection*{Eventos de herramientas}


\begin{center}
\begin{tabular}{@{}llll@{}}
\toprule
Evento & Modo & Dispatcher & Listeners principales \\
\midrule
\texttt{tools/pre-execute} & \texttt{waterfall} & \texttt{tools} & \texttt{hooks-claude-code}, \texttt{hooks-codex}, \texttt{tool-jobs} \\
\texttt{tools/execute} & \texttt{waterfall} & \texttt{tools} & \texttt{session-checkpoint-policy}, \texttt{timeout-policy} \\
\texttt{tools/post-execute} & \texttt{waterfall} & \texttt{tools} & \texttt{hooks-claude-code}, \texttt{hooks-codex}, \texttt{repeat-tool-reminder}, \texttt{spill-policy}, \texttt{tool-fs-search} \\
\texttt{tools/ptc-dispatch-log} & \texttt{waterfall} & \texttt{tools} & \texttt{spill-policy} \\
\texttt{tools/result} & \texttt{emit} & \texttt{tools} & \texttt{agent-instructions}, \texttt{subagent-in-process-driver} \\
\texttt{tools/change} & \texttt{emit} & \texttt{agent-presets}, \texttt{tools} & \texttt{tool-subagent} \\
\bottomrule
\end{tabular}
\end{center}



\subsection*{Eventos de subagente}


\begin{center}
\begin{tabular}{@{}llll@{}}
\toprule
Evento & Modo & Dispatcher & Listeners \\
\midrule
\texttt{subagent/start} & \texttt{emit} & \texttt{subagent} & \texttt{hooks-claude-code}, \texttt{subagent} \\
\texttt{subagent/end} & \texttt{emit} & \texttt{subagent} & \texttt{hooks-claude-code}, \texttt{server}, \texttt{subagent} \\
\texttt{subagent/provider-added} & \texttt{emit} & \texttt{subagent} & - \\
\bottomrule
\end{tabular}
\end{center}



\subsection*{Eventos internos no declarados}


\begin{center}
\begin{tabular}{@{}lll@{}}
\toprule
Event string & Dispatchers & Listeners \\
\midrule
\texttt{internal/dispatch} & - & \texttt{agent-team}, \texttt{commands}, \texttt{compaction}, \texttt{fs}, \texttt{goal}, \texttt{goal-round-driver}, \texttt{hook-protocol}, \texttt{llm-retry}, \texttt{permission-presets}, \texttt{plan-mode}, \texttt{sandbox-policy}, \texttt{schedule}, \texttt{scope}, \texttt{session}, \texttt{session-log-deepseek}, \texttt{session-title}, \texttt{subagent}, \texttt{terminal-bash}, \texttt{time-context}, \texttt{tool-todo}, \texttt{tool-workflow}, \texttt{tools}, \texttt{user-approval}, \texttt{webhook}, \texttt{workflow} \\
\texttt{internal/plugin} & - & \texttt{inspector}, \texttt{loader}, \texttt{lsp-stdio}, \texttt{modules} \\
\texttt{internal/service} & - & \texttt{agent-presets}, \texttt{gateway} \\
\texttt{internal/status} & - & \texttt{agent}, \texttt{inspector} \\
\bottomrule
\end{tabular}
\end{center}



\section*{configuracion}

\section*{Configuración del proyecto dsh}


\subsection*{Perfiles y arranque}


\subsubsection*{Composición de la configuración}


La configuración efectiva se construye aplicando capas en orden:


\begin{enumerate}[leftmargin=*,topsep=2pt,itemsep=1pt]
  \item Parches de los bundles listados en \texttt{dsh.profile.bundles} del manifiesto del perfil.
  \item \texttt{cordis.patch.yml} del propio perfil.
  \item \texttt{$DSH\_HOME/cordis.patch.yml} (preferencias locales de máquina, mayor prioridad que el perfil).
  \item Cada overlay \texttt{--patch <path>} en orden de aparición.
\end{enumerate}


Las capas posteriores sobrescriben filas completas (no se hace merge profundo de \texttt{config}). Se pueden insertar nuevas filas.


\subsubsection*{Recarga de parches}


\begin{itemize}[leftmargin=*,topsep=2pt,itemsep=1pt]
  \item \texttt{dsh.profile.patchReload: "live"} (por defecto en perfiles personalizados): observa cambios en archivos y reaplica.
  \item \texttt{dsh.profile.patchReload: "startup"}: carga única al iniciar.
\end{itemize}


\subsubsection*{Resolución de bundles y plugins}


\begin{itemize}[leftmargin=*,topsep=2pt,itemsep=1pt]
  \item Bundles internos (\texttt{@deepseek-ai/dsh-base}, \texttt{-web-app}, \texttt{-headless}, \texttt{-sdk-app}, \texttt{-sdk-minimal}, \texttt{-acp-app}) se resuelven desde la instalación de \texttt{dsh}.
  \item Bundles externos desde \texttt{node\_modules} del perfil (gestionado por pnpm).
  \item Plugins con nombre simple se resuelven recorriendo directorios padre hasta \texttt{$DSH\_HOME/profiles/node\_modules}.
\end{itemize}


\subsubsection*{Perfiles predefinidos}


Se autoinicializan desde plantillas en el primer uso:


\begin{center}
\begin{tabular}{@{}lll@{}}
\toprule
Perfil & Composición & patchReload \\
\midrule
\texttt{web} & base + web-app & live \\
\texttt{headless} & base + headless & startup \\
\texttt{sdk} & base + sdk-app & startup \\
\texttt{sdk-minimal} & bundle independiente & startup \\
\texttt{acp} & base + acp-app & startup \\
\bottomrule
\end{tabular}
\end{center}



Perfil faltante no predefinido → error con sugerencia: \texttt{dsh plugin --profile <name> add <package>}.


\subsubsection*{Argumentos de aplicación}


\begin{itemize}[leftmargin=*,topsep=2pt,itemsep=1pt]
  \item Flags del launcher terminan en el primer token no reconocido.
  \item El resto se pasa a \texttt{ctx.cmdlineArgs} para que los plugins de la app lo parseen.
  \item \texttt{dsh --profile web --port 8080} → \texttt{--port} va a la app web.
  \item \texttt{dsh --profile web --help} → imprime ayuda de la app y no arranca.
  \item \texttt{dsh --help} → ayuda del launcher.
  \item \texttt{-V}/\texttt{--version} antes de args de app → versión del launcher.
\end{itemize}


\subsubsection*{Precedencia de flags vs. configuración}


\begin{itemize}[leftmargin=*,topsep=2pt,itemsep=1pt]
  \item Un plugin inyecta \texttt{cmdlineArgs}, parsea argumentos y expone un servicio.
  \item Las filas de configuración usan expresiones como \texttt{port: !!js ctx.webStartup.port ?? 3080}.
  \item El flag tiene prioridad sobre el valor literal en la config.
  \item Si un parche de usuario reemplaza \texttt{config} con literales, se pierde la lectura en runtime.
  \item En modo \texttt{live}, editar un archivo de parche reevalúa expresiones pero no reinicia servicios (no puede cambiar un puerto ya servido).
\end{itemize}


\subsubsection*{Límite de argumentos}


\begin{itemize}[leftmargin=*,topsep=2pt,itemsep=1pt]
  \item \texttt{--} es consumido por el parser del launcher. Para pasar \texttt{--} literal a la app: \texttt{-- --}.
  \item Si el primer argumento de app es \texttt{web} o \texttt{plugin}, se selecciona ese subcomando.
  \item \texttt{ctx.cmdlineArgs.get()} es lectura inmutable compartida.
\end{itemize}


\subsubsection*{Apps embarcadas y sus argumentos}


\begin{center}
\begin{tabular}{@{}ll@{}}
\toprule
Perfil & Argumentos \\
\midrule
\texttt{web} & \texttt{--host}, \texttt{--port}, \texttt{--trusted-host} (repetible), \texttt{--no-open} \\
\texttt{headless} & texto de la tarea como argumento posicional \\
\texttt{sdk} & sin opciones; protocolo JSON-RPC por stdio \\
\texttt{sdk-minimal} & sin opciones; mismo protocolo JSON-RPC \\
\texttt{acp} & sin opciones; Agent Client Protocol por stdio \\
\bottomrule
\end{tabular}
\end{center}



\subsubsection*{Modo headless (one-shot)}


\begin{lstlisting}
dsh --profile headless "run the tests"
\end{lstlisting}


\begin{itemize}[leftmargin=*,topsep=2pt,itemsep=1pt]
  \item Crea un Agent persistente, envía la tarea, espera quiescencia, flushea la sesión.
  \item Último texto no vacío del asistente → stdout.
  \item Razón final \texttt{turn/end} y deltas de razonamiento no vacíos → stderr (con encabezado \texttt{dsh: reasoning:}).
  \item Exit 0 si \texttt{completed}, 1 en otro caso.
  \item Sin tarea → error de uso.
  \item No monta servidor HTTP, WebSocket, ni navegador; no abre puertos.
\end{itemize}


\subsubsection*{Volcado de configuración}


\begin{lstlisting}
dsh --profile web --dump-default-config   # solo capas de bundles
dsh --profile web --patch ./extra.yml --dump-config  # incluye parches
\end{lstlisting}


\begin{itemize}[leftmargin=*,topsep=2pt,itemsep=1pt]
  \item Imprime comentarios con archivo de origen y overlays que modificaron cada fila.
  \item Expresiones \texttt{!!js} sin evaluar.
  \item Nombres de plugin relativos se resuelven junto al archivo de parche.
  \item Destinos de parche no encontrados → stderr.
  \item No ejecuta providers de línea de comandos de la app; rechaza si hay argumentos de app.
\end{itemize}


\subsection*{Gestión de plugins (\texttt{dsh plugin})}


\begin{lstlisting}
dsh plugin --profile <name> <args...>
\end{lstlisting}


\begin{itemize}[leftmargin=*,topsep=2pt,itemsep=1pt]
  \item Si el perfil no existe, lo inicializa (plantilla predefinida o solo \texttt{@deepseek-ai/dsh-base}).
  \item Reenvía \texttt{<args...>} a \texttt{pnpm} con el directorio del perfil como working directory.
  \item pnpm debe estar en el PATH.
  \item Especificaciones de ruta relativas (\texttt{.}, \texttt{../plugin}, \texttt{file:}, \texttt{link:}) se anclan al directorio de invocación.
  \item Tras cada ejecución exitosa, reconcilia \texttt{dsh.profile.bundles}:
  \item Dependencia con \texttt{"dsh": { "bundle": { "patch": "./cordis.patch.yml" } }} → se añade a la pila de capas.
  \item Dependencia sin declaración de bundle → warning único.
  \item Dependencia eliminada → se quita de la pila.
\end{itemize}


\subsubsection*{Subagents Codex y Claude Code (bundles opcionales)}


\begin{lstlisting}
dsh plugin --profile <name> add @deepseek-ai/dsh-subagent-codex
dsh plugin --profile <name> add @deepseek-ai/dsh-subagent-claude-code
dsh plugin --profile <name> add @deepseek-ai/dsh-subagent-codex @deepseek-ai/dsh-subagent-claude-code
dsh plugin --profile <name> remove @deepseek-ai/dsh-subagent-codex
dsh plugin --profile <name> remove @deepseek-ai/dsh-subagent-claude-code
\end{lstlisting}


\begin{itemize}[leftmargin=*,topsep=2pt,itemsep=1pt]
  \item La operación de pnpm cambia manifiesto y lista de bundles en disco.
  \item Un perfil en ejecución mantiene el conjunto de bundles de su inicio actual.
  \item \textbf{Reiniciar el perfil} después de añadir/quitar/actualizar bundles.
  \item Ediciones en \texttt{cordis.patch.yml} del perfil o home se aplican por hot reload.
\end{itemize}


\subsubsection*{Plugins desde Git}


\begin{lstlisting}
dsh plugin --profile tui add github:deepseek-harness/turtle-ui
dsh plugin --profile tui remove turtle-ui
dsh --profile tui
\end{lstlisting}


\begin{itemize}[leftmargin=*,topsep=2pt,itemsep=1pt]
  \item Plugins con fuentes necesitan \texttt{prepare} script → pnpm ≥10 bloquea hasta permitir en \texttt{pnpm-workspace.yaml}.
  \item Primera ejecución de \texttt{add} falla con hint; copiar la clave impresa en \texttt{pnpm-workspace.yaml} y re-ejecutar.
  \item Tarballs preconstruidos o checkouts locales no requieren permisos.
\end{itemize}


\subsection*{Alias web (\texttt{dsh web})}


Alias fijo para \texttt{--profile web}. Los flags después de \texttt{web} pertenecen a la app web.


\begin{lstlisting}
dsh web
dsh web --no-open
dsh web --patch ./extra.cordis.yml
dsh web --dump-config
dsh web --help
\end{lstlisting}


\subsubsection*{Flags de la app web}


\begin{itemize}[leftmargin=*,topsep=2pt,itemsep=1pt]
  \item \texttt{--host}, \texttt{--port}: sobrescriben valores compuestos.
  \item \texttt{--trusted-host} (repetible): añade autoridades a \texttt{ctx.webRuntime.trustedHosts}.
  \item \texttt{--no-open}: deshabilita apertura de navegador.
\end{itemize}


\subsubsection*{Comportamiento del servidor}


\begin{itemize}[leftmargin=*,topsep=2pt,itemsep=1pt]
  \item Por defecto sirve en \texttt{http://127.0.0.1:3080}.
  \item Solo abre navegador tras asentarse el árbol Loader completo.
  \item \texttt{SSH\_CONNECTION} o \texttt{SSH\_TTY} no vacíos → suprime handoff del navegador (imprime URL).
  \item No soporta \texttt{--host 0.0.0.0} (error de uso).
  \item HMR de cliente siempre montado, inactivo hasta que \texttt{pnpm run dev:web} reconstruya bundles.
\end{itemize}


\subsection*{cordis.yml}


\begin{lstlisting}
# cordis.yml: lista de entradas. Cada entrada acepta:
- id: identidad estable para diff del loader
- name: ruta al módulo o nombre de paquete
- disabled: true  # conserva la entrada pero no monta el plugin
- config: bloque validado por el schema del plugin
- !!js: solo dentro de config y disabled; evalúa expresión JS en tiempo de carga
\end{lstlisting}


\begin{itemize}[leftmargin=*,topsep=2pt,itemsep=1pt]
  \item \textbf{Grupos}: anidan sub-listas de entradas; \texttt{isolate} da a un grupo su propia instancia de un servicio.
  \item \textbf{HMR}: \texttt{@deepseek-ai/cordis-plugin-hmr} vigila archivos y recarga plugins al guardar. Requiere \texttt{@deepseek-ai/cordis-plugin-logger-console} y \texttt{@deepseek-ai/cordis-plugin-timer} como dependencias.
\end{itemize}


\subsection*{Interfaces de configuración de plugins}


\subsubsection*{\texttt{@deepseek-ai/dsh-settings-file}}


\begin{lstlisting}
interface Config {
  path?: string;       // Ruta al archivo. Default: `settings.yaml` en `dshHome`.
  dshHome?: string;    // Directorio home del harness. Default: `$DSH_HOME` o `~/.dsh`.
  watch?: boolean;     // Observar cambios y recargar. Default: `true`.
  debounceMs?: number; // Ventana de estabilización del watcher (ms). Default: `100`.
}
\end{lstlisting}


\subsubsection*{\texttt{@deepseek-ai/dsh-shell-env}}


\begin{lstlisting}
interface Config {
  dshHome?: string; // Default: `$DSH_HOME` o `~/.dsh`.
}
\end{lstlisting}


\subsubsection*{\texttt{@deepseek-ai/dsh-skill}}


\begin{lstlisting}
interface Config {
  collectCacheMaxEntries?: number; // Máximo de catálogos completados en memoria.
}
\end{lstlisting}


\subsubsection*{\texttt{@deepseek-ai/dsh-skill-filesystem}}


\textbf{Requiere:} \texttt{skills}


\begin{lstlisting}
interface Config {
  providerName?: string;            // Default: `filesystem`.
  includeDefaultRoots?: boolean;
  dshHome?: string;                 // Default: `$DSH_HOME` o `~/.dsh`.
  agentsHome?: string;              // Default: `$DSH_AGENTS_HOME` o `~/.agents`.
  customSkillDirs?: string[];
  watch?: boolean;
  watchUsePolling?: boolean;
  watchStabilityThresholdMs?: number;
  watchPollIntervalMs?: number;
  watchMaxProjects?: number;
  watchFollowSymlinks?: boolean;
  bundledSkillDir?: string;         // Default: `$DSH_BUNDLED_SKILL_DIR`.
}
\end{lstlisting}


\subsubsection*{\texttt{@deepseek-ai/dsh-spill-local}}


\begin{lstlisting}
interface Config {
  root?: string;             // Directorio raíz. Default: directorio temporal privado (0700).
  cleanupPeriodDays?: number; // Días para limpieza de archivos viejos. Default: `30`. `0` desactiva.
}
\end{lstlisting}


\subsubsection*{\texttt{@deepseek-ai/dsh-spill-policy}}


\textbf{Requiere:} \texttt{tools}


\begin{lstlisting}
interface Config {
  maxInlineBytes?: number; // Bytes máx. en texto plano antes de spill. Omitir desactiva la política.
}
\end{lstlisting}


\subsubsection*{\texttt{@deepseek-ai/dsh-storage-domain}}


\textbf{Requiere:} \texttt{storage}


\begin{lstlisting}
interface Config {
  backend: string;                  // Backend por defecto (requerido).
  routes?: Record<string, string>;  // Overrides por dominio: nombre -> backend.
}
\end{lstlisting}


\subsubsection*{\texttt{@deepseek-ai/dsh-storage-json}}


\textbf{Requiere:} \texttt{storage}


\begin{lstlisting}
interface Config {
  root: string; // Directorio con archivos `<unit>.json` (requerido, sin default).
}
\end{lstlisting}


\subsubsection*{\texttt{@deepseek-ai/dsh-storage-sqlite}}


\textbf{Requiere:} \texttt{storage}


\begin{lstlisting}
type JournalMode = 'wal' | 'delete' | 'truncate' | 'persist';

interface Config {
  path: string;                // Ruta al archivo de BD. Usar `:memory:` para tests.
  journalMode?: JournalMode;   // Default: `wal`.
}
\end{lstlisting}


\subsubsection*{\texttt{@deepseek-ai/dsh-subagent-acp}}


\textbf{Requiere:} \texttt{subagents}, \texttt{subprocess}


\begin{lstlisting}
type PermissionPolicy = 'allow' | 'reject';

interface Config {
  providerName: string;           // Default: `acp`.
  command: string;                // Ejecutable del proceso hijo.
  args: string[];                 // Argumentos para el comando.
  cwd?: string;                   // Directorio de trabajo. Default: hereda del padre.
  permission: PermissionPolicy;   // `reject` (default) o `allow`.
  env: Record<string, string>;    // Variables de entorno extra.
  disposeEofGraceMs?: number;
  disposeGraceMs?: number;
}
\end{lstlisting}


\subsubsection*{\texttt{@deepseek-ai/dsh-subagent-claude-code}}


\textbf{Requiere:} \texttt{subagents}, \texttt{subprocess}


\begin{lstlisting}
interface Config {
  providerName?: string;          // Default: `claude-code`.
  model?: string;
  env?: Record<string, string>;
  permissionMode?: ClaudeCodePermissionMode;
  disposeGraceMs?: number;
}
\end{lstlisting}


\subsubsection*{\texttt{@deepseek-ai/dsh-subagent-codex}}


\textbf{Requiere:} \texttt{subagents}, \texttt{subprocess}


\begin{lstlisting}
type CodexPermissionMode = 'never' | 'approve-for-me' | 'dangerously-bypass-approvals-and-sandbox';

interface Config {
  providerName?: string;  // Default: `codex`.
  model?: string;
  env?: Record<string, string>;
  permissionMode?: CodexPermissionMode;
  disposeGraceMs?: number;
}
\end{lstlisting}


\subsubsection*{\texttt{@deepseek-ai/dsh-subagent-dsh-sdk}}


\textbf{Requiere:} \texttt{subagents}


\begin{lstlisting}
interface Config {
  providerName: string;       // Default: `dsh-sdk`.
  dshBin?: string;
  profile: string;            // Perfil del hijo (default: `sdk`).
  patches: string[];
  dshHome: string;            // Home aislado para el hijo.
  cwd?: string;
  provider: string;           // Default: `deepseek-official`.
  model: string;              // Default: `deepseek-v4-flash`.
  maxTokens?: number;
  env: Record<string, string>;
  shutdownTimeoutMs?: number;
  disposeEofGraceMs?: number;
  disposeGraceMs?: number;
}
\end{lstlisting}


\subsubsection*{\texttt{@deepseek-ai/dsh-subagent-fork-in-process}}


\textbf{Requiere:} \texttt{subagents}


\begin{lstlisting}
interface Config {
  providerName: string; // Default: `fork`.
}
\end{lstlisting}


\subsubsection*{\texttt{@deepseek-ai/dsh-subagent-spawn-in-process}}


\textbf{Requiere:} \texttt{subagents}


\begin{lstlisting}
interface Config {
  providerName: string; // Default: `spawn`.
}
\end{lstlisting}


\subsubsection*{\texttt{@deepseek-ai/dsh-subprocess-e2b}}


\textbf{Requiere:} \texttt{e2b}


\begin{lstlisting}
interface Config {
  pollMs?: number; // Cadencia de sondeo de estado (ms).
}
\end{lstlisting}


\subsubsection*{\texttt{@deepseek-ai/dsh-system-prompt}}


\begin{lstlisting}
interface Config {
  includeHarnessIdentity?: boolean; // Default: `true`.
  includeRuntimeContext?: boolean;  // Default: `true`.
  personaPrefix?: string;           // Plantilla de prefijo de persona.
  personaSuffix?: string;           // Plantilla de sufijo de persona.
  toolOrder?: string[];             // Orden de herramientas. Default: lexicográfico.
}
\end{lstlisting}


\subsubsection*{\texttt{@deepseek-ai/dsh-terminal-bash}}


\textbf{Requiere:} \texttt{terminals}, \texttt{sandboxPolicy}, \texttt{sessionProjections}, \texttt{subprocess}


\begin{lstlisting}
type ShellDialect = 'bash' | 'pwsh';

interface Config {
  backendType?: string;       // Default: `shell`.
  shellDialect?: ShellDialect; // Default: `bash`.
  shellPath?: string;         // Default: `/bin/bash`.
  shellArgs?: string[];       // Default: `--noprofile --norc -i`.
  rows?: number;
  cols?: number;
  scrollbackLines?: number;
  scrollbackMaxBytes?: number;
  maxReadBytes?: number;
  pollIntervalMs?: number;
  exactProbeAfterMs?: number;
  idleSilenceMs?: number;
  handoffGraceMs?: number;
  timeoutMs?: number;
  disposeGraceMs?: number;
}
\end{lstlisting}


\subsubsection*{\texttt{@deepseek-ai/dsh-time-context}}


\textbf{Requiere:} \texttt{agents}, \texttt{sessionProjections}


\begin{lstlisting}
interface Config {
  timeZone?: string;            // Zona horaria. Default: zona del proceso.
  refreshIntervalMs?: number;   // Mínimo ms entre inyecciones. `0` = cada paso.
}
\end{lstlisting}


\subsubsection*{\texttt{@deepseek-ai/dsh-typert-loader}}


\begin{lstlisting}
interface Config {
  packages?: string[]; // Paquetes npm adicionales que deben resolver y exportar `./typert`.
}
\end{lstlisting}


\subsubsection*{\texttt{@deepseek-ai/dsh-user-approval}}


\begin{lstlisting}
type ApprovalPolicy = 'ask' | 'never'

interface Config {
  readonly policy?: ApprovalPolicy
  // 'ask' delega en los answerers compuestos (fail-closed si no hay).
  // 'never' rechaza automáticamente sin preguntar (CI/unattended).
}
\end{lstlisting}


\subsubsection*{\texttt{@deepseek-ai/dsh-web}}


\begin{lstlisting}
interface WebRuntimeConfig {
  readonly searchProvider?: string  // ID explícito del proveedor de búsqueda.
  readonly fetchProvider?: string   // ID explícito del proveedor de fetch.
}
\end{lstlisting}


\subsubsection*{\texttt{@deepseek-ai/dsh-web-app}}


\begin{lstlisting}
interface Config {
  openBrowser: boolean
  printUrl: boolean
  surfaceContext: boolean
  trustedHosts: string[]
}
\end{lstlisting}


\subsubsection*{\texttt{@deepseek-ai/dsh-web-fetch-http}}


\begin{lstlisting}
interface Config {
  maxResponseBytes?: number
  maxBodyChars?: number
  timeoutMs?: number
  maxRedirects?: number
  userAgent?: string
}
\end{lstlisting}


\subsubsection*{\texttt{@deepseek-ai/dsh-web-search-deepseek}}


\begin{lstlisting}
interface Config {
  apiKey?: string
  apiKeyEnv?: string       // Default: `DEEPSEEK_API_KEY`
  baseURL?: string
  model?: string           // Default: `deepseek-v4-flash`
  apiVersion?: string      // Default: `2023-06-01`
  maxTokens?: number       // Default: 4096
  maxUses?: number         // Default: 5
}
\end{lstlisting}


\subsubsection*{\texttt{@deepseek-ai/dsh-web-search-exa}}


\begin{lstlisting}
interface Config {
  apiKey?: string
  baseURL?: string
  searchType?: 'auto' | 'keyword' | 'neural'  // Default: `auto`
  numResults?: number
  highlightsPerResult?: number  // Default: 1
}
\end{lstlisting}


\subsubsection*{\texttt{@deepseek-ai/dsh-web-search-perplexity}}


\begin{lstlisting}
interface Config {
  apiKey?: string
  baseURL?: string
  model?: string           // Default: `sonar`
  maxTokens?: number       // Default: 1024
  searchRecency?: 'day' | 'week' | 'month' | 'year'
}
\end{lstlisting}


\subsubsection*{\texttt{@deepseek-ai/dsh-webhook-github}}


\begin{lstlisting}
interface Config {
  readonly source: string
  readonly path: string
  readonly secretEnv: string
  readonly maxBodyBytes: number
}
\end{lstlisting}


\subsubsection*{\texttt{@deepseek-ai/dsh-workflow-worker-thread}}


\begin{lstlisting}
interface Config {
  provider?: string              // Default: `spawn`
  maxConcurrentAgents?: number   // `0` auto-resuelve a `min(16, max(1, cores - 2))`
  maxTotalAgents?: number        // Default: 1000
  maxItemsPerCall?: number       // Default: 4096
  syncTimeoutMs?: number         // Default: 5000 ms
  disposeGraceMs?: number        // Default: 5000 ms
}
\end{lstlisting}


\subsubsection*{\texttt{@deepseek-ai/dsh-experimental-agent-team}}


\begin{lstlisting}
interface Config {
  readonly maxMembers?: number
  readonly maxTasks?: number
  readonly maxPendingMessagesPerMember?: number
  readonly maxMessageBytes?: number
  readonly disposalTimeoutMs?: number
}
\end{lstlisting}


\subsubsection*{\texttt{@deepseek-ai/dsh-experimental-code-runtime-python}}


\begin{lstlisting}
interface Config {
  cpuSeconds?: number
  maxWallMs?: number
  addressSpaceMb?: number
  maxLogBytes?: number
  maxValueBytes?: number
  graceMs?: number
  pythonBin?: string
}
\end{lstlisting}


\subsubsection*{\texttt{@deepseek-ai/dsh-experimental-inspector}}


\begin{lstlisting}
interface Config extends Omit<InspectorOptions, 'clientOrigins'> {
  clientOrigins?: string[]
}

interface InspectorOptions {
  readonly host?: '127.0.0.1'
  readonly port?: number
  readonly clientOrigins?: readonly string[]
  readonly captureFetch?: boolean
  readonly maxRequestBodyBytes?: number
  readonly maxResponseBodyBytes?: number
  readonly maxBodyChunkBytes?: number
  readonly maxJournalBytes?: number
  readonly maxRetainedRequests?: number
  readonly maxSourceFrameBytes?: number
  readonly maxSourceRecordsPerFrame?: number
  readonly maxQueuedRecords?: number
  readonly maxQueuedBytes?: number
  readonly startupTimeoutMs?: number
  readonly stopTimeoutMs?: number
  readonly clientReconnectBaseMs?: number
  readonly clientReconnectMaxMs?: number
  readonly clientRuntimeTimeoutMs?: number
  readonly queryTimeoutMs?: number
  readonly maxClientRuntimeObjects?: number
  readonly maxClientRuntimeProperties?: number
  readonly maxClientSourceBytes?: number
  readonly maxCordisNodes?: number
  readonly maxDisconnectedCordisTrees?: number
}
\end{lstlisting}


\subsubsection*{\texttt{@deepseek-ai/dsh-experimental-tool-agent-team}}


\begin{lstlisting}
interface Config {
  readonly freshProvider?: string
  readonly forkProvider?: string
}
\end{lstlisting}


\subsubsection*{\texttt{@deepseek-ai/dsh-file-reference-local}}


\begin{lstlisting}
interface Config {
  maxResults?: number
  maxEntries?: number
  excludedDirectories?: string[]
}
\end{lstlisting}


\subsubsection*{\texttt{@deepseek-ai/dsh-fs-local}}


\begin{lstlisting}
interface Config {
  cwd?: string
  diffBasisMaxBytes?: number  // Default: 10 MiB
}
\end{lstlisting}


\subsubsection*{\texttt{@deepseek-ai/dsh-fs-sandbox}}


\begin{lstlisting}
type Config = LocalConfig; // Misma configuración que fs-local
\end{lstlisting}


\subsubsection*{\texttt{@deepseek-ai/dsh-goal}}


\begin{lstlisting}
interface Config {
  defaultMaxGoalRounds?: number
}
\end{lstlisting}


\subsubsection*{\texttt{@deepseek-ai/dsh-headless}}


\begin{lstlisting}
interface Config {
  task: string
}
\end{lstlisting}


\subsubsection*{\texttt{@deepseek-ai/dsh-hooks-claude-code}}


\begin{lstlisting}
interface Config {
  configPath: string
  pluginRoot?: string
  projectDir?: string
  defaultTimeoutMs?: number        // Default: 600000
  stderrSummaryMaxChars?: number
}
\end{lstlisting}


\subsubsection*{\texttt{@deepseek-ai/dsh-hooks-codex}}


\begin{lstlisting}
interface Config {
  configPath: string
  model?: string
  defaultTimeoutMs?: number        // Default: 600000
  stderrSummaryMaxChars?: number
}
\end{lstlisting}


\subsubsection*{\texttt{@deepseek-ai/dsh-host-directory-picker-browse}}


\begin{lstlisting}
interface Config {
  maxEntries: number
}
\end{lstlisting}


\subsubsection*{\texttt{@deepseek-ai/dsh-host-frontend-static}}


\begin{lstlisting}
interface Config {
  distIndex: string
}
\end{lstlisting}


\subsubsection*{\texttt{@deepseek-ai/dsh-host-open-in-app}}


\begin{lstlisting}
interface Config {
  readonly probeTimeoutMs: number
  readonly iconTimeoutMs: number
  readonly launchWatchMs: number
}
\end{lstlisting}


\subsubsection*{\texttt{@deepseek-ai/dsh-host-webserver}}


\begin{lstlisting}
interface Config {
  host: '127.0.0.1' | '0.0.0.0'
  port: number
  compression?: 'none' | 'gzip'    // Default: 'none'
  compressionLevel?: number        // Default: 1
  compressionThresholdBytes?: number // Default: 1024
}
\end{lstlisting}


\subsubsection*{\texttt{@deepseek-ai/dsh-invariants}}


\begin{lstlisting}
interface Config {
  readonly enabled?: boolean           // Default: true
  readonly package_allowlist?: string[]
  readonly package_blocklist?: string[]
}
\end{lstlisting}


\subsubsection*{\texttt{@deepseek-ai/dsh-jobs-local}}


\begin{lstlisting}
interface Config {
  maxConcurrentJobsPerOwner?: number  // Default: 10
}
\end{lstlisting}


\subsubsection*{\texttt{@deepseek-ai/dsh-agent-loop}}


\begin{lstlisting}
interface Config {
  maxParallelToolCalls?: number
  agents?: Array<{
    id: string
    provider: string
    model: string
    cwd?: string
    sessionId?: string
    resumeSessionId?: string
  }>
}
\end{lstlisting}


\subsubsection*{\texttt{@deepseek-ai/dsh-compaction-basic}}


\begin{lstlisting}
interface Config {
  thresholdRatio?: number
  retainRatio?: number
  summarizationProvider?: string
  summarizationModel?: string
  maxTokens?: number
}
\end{lstlisting}


\subsubsection*{\texttt{@deepseek-ai/dsh-bash-local}}


\begin{lstlisting}
interface Config {
  timeoutMs?: number
  maxOutputBytes?: number
  graceMs?: number
}
\end{lstlisting}


\subsubsection*{\texttt{@deepseek-ai/dsh-code-runtime-worker-thread}}


\begin{lstlisting}
interface Config {
  computeMs?: number
  maxWallMs?: number
  maxOldGenerationSizeMb?: number
}
\end{lstlisting}


\subsubsection*{\texttt{@deepseek-ai/dsh-attachment-local}}


\begin{lstlisting}
interface Config {
  maxImageBytes?: number
  maxImagePixels?: number
  normalizedImageMaxBytes?: number
}
\end{lstlisting}


\subsubsection*{\texttt{@deepseek-ai/dsh-credentials-local}}


\begin{lstlisting}
interface Config {
  path?: string
  dshHome?: string
  watch?: boolean
  debounceMs?: number
}
\end{lstlisting}


\subsection*{Interfaces de configuración adicionales}


\subsubsection*{\texttt{ReplayProviderConfig}}


\begin{lstlisting}
interface ReplayProviderConfig {
  id: string
  name?: string
  models?: ReplayModelConfig[]
  retryPolicy?: RetryPolicyConfig
}
\end{lstlisting}


\subsubsection*{\texttt{ReplayModelConfig}}


\begin{lstlisting}
interface ReplayModelConfig {
  id: string
  name?: string
  description?: string
  contextWindow?: number
  inputModalities?: readonly ModelModality[]
  defaultMaxTokens?: number
  imageRequestTokens?: number
  reasoningEfforts?: string[]
  defaultReasoningEffort?: string
}
\end{lstlisting}


\subsubsection*{\texttt{LspLocalServerConfig}}


\begin{lstlisting}
interface LspLocalServerConfig {
  command: string
  extensionToLanguage: Record<string, string>
  args?: string[]
  env?: Record<string, string>
  initializationOptions?: unknown
  configuration?: unknown
  maxMessageBytes?: number
  maxStderrBytes?: number
  maxDocumentBytes?: number
  shutdownTimeoutMs?: number
  killGraceMs?: number
}
\end{lstlisting}


\subsubsection*{\texttt{StdioConfig}}


\begin{lstlisting}
interface StdioConfig {
  transport: 'stdio'
  serverName: string
  command: string
  args: string[]
  env: Record<string, string>
  cwd: string
  toolCallTimeoutMs: number
  failOnStartupError: boolean
  reconnect?: ReconnectConfig
}
\end{lstlisting}


\subsubsection*{\texttt{StreamableHttpConfig}}


\begin{lstlisting}
interface StreamableHttpConfig {
  transport: 'streamable-http'
  serverName: string
  url: string
  headers: Record<string, string>
  toolCallTimeoutMs: number
  failOnStartupError: boolean
  reconnect?: ReconnectConfig
}
\end{lstlisting}


\subsubsection*{\texttt{ReconnectConfig}}


\begin{lstlisting}
interface ReconnectConfig {
  enabled?: boolean
  initialDelayMs?: number
  maxDelayMs?: number
  maxAttempts?: number
}
\end{lstlisting}


\subsubsection*{\texttt{PresetSpec}}


\begin{lstlisting}
interface PresetSpec {
  sandbox: SandboxMode
  approval: ApprovalPolicy
  name?: string
  description?: string
}
\end{lstlisting}


\subsubsection*{\texttt{PlanModeConfig}}


\begin{lstlisting}
interface PlanModeConfig {
  section: string
}
\end{lstlisting}


\subsubsection*{\texttt{JsonRpcConfig}}


\begin{lstlisting}
interface JsonRpcConfig {
  maxTokensAsSuccess?: boolean
  input?: Readable
  output?: Writable
  exit?: (code: number) => void
}
\end{lstlisting}


\subsubsection*{\texttt{SessionQueryConfig}}


\begin{lstlisting}
interface Config extends SessionQueryConfig {
  path: string
  openAt?: OpenAt
  journalMode?: JournalMode
  defaultLimit?: number
  maxLimit?: number
  snippetChars?: number
  persistedReadConcurrency?: number
  preparedSessionCacheSize?: number
}
\end{lstlisting}


\subsubsection*{\texttt{SessionTelemetryMode}}


\begin{lstlisting}
enum SessionTelemetryMode {
  FEEDBACK_ONLY = 'FEEDBACK_ONLY',
  DISABLED = 'DISABLED',
}
\end{lstlisting}


\subsubsection*{\texttt{ToolSubagentConfig}}


\begin{lstlisting}
interface Config {
  provider: string
  toolName?: string
  modelSelectionSettings?: boolean
  enableRunInBackground?: boolean
  backgroundMode?: 'one-shot' | 'continuable'
  agentOptions?: AgentOptions
  persona?: string
  toolFilter?: { allow?: string[]; deny?: string[] }
  maxDepth?: number | 'provider-managed'
}
\end{lstlisting}


\subsubsection*{\texttt{ToolPresentationMode}}


\begin{lstlisting}
type ToolPresentationMode = 'native' | 'ptc' | 'both'
\end{lstlisting}


\subsection*{Variables de entorno}


\begin{itemize}[leftmargin=*,topsep=2pt,itemsep=1pt]
  \item \texttt{$DSH\_HOME}: directorio home del harness. Default: \texttt{\textasciitilde{}/.dsh}.
  \item \texttt{$DSH\_AGENTS\_HOME}: directorio de agentes. Default: \texttt{\textasciitilde{}/.agents}.
  \item \texttt{$DSH\_BUNDLED\_SKILL\_DIR}: directorio de skills empaquetados.
  \item \texttt{SSH\_CONNECTION} / \texttt{SSH\_TTY}: si no están vacíos, suprime handoff del navegador en modo web.
  \item \texttt{DEEPSEEK\_API\_KEY}: clave de API para el proveedor DeepSeek.
  \item \texttt{CLAUDE\_PLUGIN\_ROOT}: reemplazado en cadenas de comando de hooks Claude Code.
  \item \texttt{CLAUDE\_PROJECT\_DIR}: reemplazado en cadenas de comando de hooks Claude Code.
\end{itemize}


\subsection*{Comandos de desarrollo}


\begin{lstlisting}
pnpm dsh web              # Inicia el host desde fuente (usa el fallback SRC)
pnpm run dev:web          # Inicia el watcher de plugins del cliente
pnpm run build:lib        # Reconstruye los contratos estrictos de Typert (Host y Cliente)
pnpm run build:lib:host   # Genera artefactos estrictos del Host
pnpm run build:lib:client # Genera artefactos estrictos del Cliente
\end{lstlisting}


\section*{ejemplos}

\section*{Ejemplos de código real del repositorio}


\subsection*{Ejemplo 1: Declaración de un servicio remoto (Host)}


\textbf{Fuente}: patrón de \texttt{TypertRemoteService} en el repo


\begin{lstlisting}
import type { Agent } from '@deepseek-ai/dsh-agent'
import { TypertRemoteService, Remote, RemoteScope } from '@deepseek-ai/dsh-typert-protocol'
import type { Context } from '@deepseek-ai/cordis'

export class GoalService extends TypertRemoteService {
  constructor(ctx: Context) {
    super(ctx, 'goals')
  }

  @Remote('create')
  createForClient(
    agent: Agent,
    request: { objective: string },
    signal: AbortSignal,
  ): { accepted: boolean } {
    signal.throwIfAborted()
    return this.create(agent, request)
  }

  @RemoteScope('agent', 'current')
  currentForClient(): { accepted: boolean } {
    return { accepted: true }
  }

  private create(_agent: Agent, request: { objective: string }): { accepted: boolean } {
    return { accepted: request.objective.length > 0 }
  }
}
\end{lstlisting}


\subsection*{Ejemplo 2: Consumo de un servicio remoto (Cliente)}


\textbf{Fuente}: patrón de consumo remoto en el repo


\begin{lstlisting}
import type { Context } from '@deepseek-ai/cordis'
import type {} from '@deepseek-ai/dsh-api-remotes/client'

export const inject = ['remote', 'remote.goals']

declare const ctx: Context
declare const agentId: string

// Llamada directa
await ctx.remote.goals.create(agentId, { objective: 'ship it' })

// Llamada con ámbito de agente
declare const agentCtx: any
await agentCtx.remote.goals.create({ objective: 'ship it' })
\end{lstlisting}


\subsection*{Ejemplo 3: Llamada remota con manejo de errores}


\textbf{Fuente}: patrón de manejo de errores remotos en el repo


\begin{lstlisting}
import type { Context } from '@deepseek-ai/cordis'
import { isRemoteFailure } from '@deepseek-ai/dsh-api-gateway/client'
import type {} from '@deepseek-ai/dsh-api-remotes/client'

export const inject = ['remote', 'remote.notes']

declare const ctx: Context

export async function noteTitles(): Promise<readonly string[]> {
  const result = await ctx.remote.notes.list()
  if (!result.ok) {
    if (result.error.code === 'note/not-found') return []
    throw result.error
  }
  return result.value.map(row => row.title)
}

export async function renderTitles(): Promise<string> {
  try {
    return (await noteTitles()).join(', ')
  } catch (error: unknown) {
    if (!isRemoteFailure(error)) throw error
    return `unavailable (${error.code})`
  }
}

export function hostLabel(): string {
  const { home, isLoopback } = ctx.remote.$host
  return home ?? (isLoopback ? 'local host' : 'remote host')
}
\end{lstlisting}


\subsection*{Ejemplo 4: Prueba de Owner con remoteErrorOf}


\textbf{Fuente}: patrón de pruebas de owner en el repo


\begin{lstlisting}
import { remoteErrorOf } from '@deepseek-ai/dsh-typert-protocol'
import { expect, it } from 'vitest'

declare function rename(noteId: string, title: string): Promise<void>

it('refuses an unknown note before writing', async () => {
  const failure = await rename('n-404', 'fresh title').catch((error: unknown) => error)
  expect(remoteErrorOf(failure)).toMatchObject({
    code: 'note/not-found',
    details: { noteId: 'n-404' },
  })
})
\end{lstlisting}


\subsection*{Ejemplo 5: Prueba de Cliente con TestRemote}


\textbf{Fuente}: patrón de pruebas de cliente en el repo


\begin{lstlisting}
import { Context } from '@deepseek-ai/cordis'
import { RemoteError, TestRemote } from '@deepseek-ai/dsh-client-test-runtime'
import { expect, it } from 'vitest'

it('renders the failure code the Host reported', async () => {
  const ctx = new Context()
  const remote = new TestRemote(ctx, {
    notes: {
      list: () => Promise.resolve({
        ok: false as const,
        error: new RemoteError('note/not-found', 'no note "n-404"', { noteId: 'n-404' }),
      }),
    },
  })
  remote.$host = { home: '/home/fixture', isLoopback: true }
  await expect(ctx.remote.notes.list()).resolves.toMatchObject({ error: { code: 'note/not-found' } })
})
\end{lstlisting}


\subsection*{Ejemplo 6: Herramienta del harness}


\textbf{Fuente}: patrón de herramienta en el repo


\begin{lstlisting}
// greet-tool.ts
import type { Context } from '@deepseek-ai/cordis'
import { brandString } from '@deepseek-ai/dsh-brand'
import { defineTool } from '@deepseek-ai/dsh-tools'
import type { ToolCallId } from '@deepseek-ai/dsh-llm'

export const name = 'greet-tool'
export const inject = ['tools']

export function apply(ctx: Context) {
  ctx.tools.register(defineTool({
    name: 'greet',
    description: 'Greet the named person.',
    parameters: {
      name: { type: 'string', required: true, description: 'Who to greet' },
    },
    output: {
      schema: { type: 'string' },
      render: (_args, value) => [{ type: 'text', text: value }],
    },
    async execute(args) {
      return `Hello, ${args.name}!`
    },
  }))

  void (async () => {
    const result = await ctx.tools.execute({
      callId: brandString<ToolCallId>('demo-1'),
      name: 'greet',
      arguments: { name: 'Cordis' },
      signal: new AbortController().signal,
    })
    console.log('tool replied:', JSON.stringify(result.content))
  })()
}
\end{lstlisting}


\begin{lstlisting}
# cordis.yml
- name: '@deepseek-ai/dsh-system-prompt'
- name: '@deepseek-ai/dsh-tools'
- name: './tool-logger.ts'
- name: './greet-tool.ts'
\end{lstlisting}


\subsection*{Ejemplo 7: Configuración de agente en cordis.yml}


\textbf{Fuente}: configuración de \texttt{@deepseek-ai/dsh-agent-loop}


\begin{lstlisting}
plugins:
  '@deepseek-ai/dsh-agent-loop':
    config:
      maxParallelToolCalls: 4
      agents:
        - id: main-agent
          provider: deepseek
          model: deepseek-chat
          cwd: /home/user/project
\end{lstlisting}


\subsection*{Ejemplo 8: Configuración de compactación básica}


\textbf{Fuente}: configuración de \texttt{@deepseek-ai/dsh-compaction-basic}


\begin{lstlisting}
plugins:
  '@deepseek-ai/dsh-compaction-basic':
    config:
      thresholdRatio: 0.8
      retainRatio: 0.2
      summarizationProvider: deepseek
      summarizationModel: deepseek-chat
      maxTokens: 4096
\end{lstlisting}


\subsection*{Ejemplo 9: Servidor MCP por stdio}


\textbf{Fuente}: configuración de \texttt{@deepseek-ai/dsh-mcp-client}


\begin{lstlisting}
plugins:
  '@deepseek-ai/dsh-mcp-client':
    transport: stdio
    serverName: my-server
    command: node
    args: ['server.js']
    env: {}
    cwd: /app
    toolCallTimeoutMs: 30000
    failOnStartupError: true
    reconnect:
      enabled: true
      initialDelayMs: 500
      maxDelayMs: 30000
      maxAttempts: 10
\end{lstlisting}


\subsection*{Ejemplo 10: Subagente ACP}


\textbf{Fuente}: configuración de \texttt{@deepseek-ai/dsh-subagent-acp}


\begin{lstlisting}
plugins:
  '@deepseek-ai/dsh-subagent-acp':
    providerName: acp
    command: /usr/local/bin/acp-agent
    args: ['--config', 'agent.yaml']
    permission: allow
    env:
      DEEPSEEK_API_KEY: sk-...
    disposeEofGraceMs: 5000
    disposeGraceMs: 2000
\end{lstlisting}


\subsection*{Ejemplo 11: Herramienta subagente con filtro}


\textbf{Fuente}: configuración de \texttt{@deepseek-ai/dsh-tool-subagent}


\begin{lstlisting}
plugins:
  '@deepseek-ai/dsh-tool-subagent':
    provider: spawn
    toolName: my-subagent
    enableRunInBackground: true
    backgroundMode: continuable
    agentOptions:
      model: deepseek-v4-flash
    toolFilter:
      allow: ['read_file', 'write_file']
      deny: ['shell']
    maxDepth: 2
\end{lstlisting}


\subsection*{Ejemplo 12: Presentación de herramientas en modo PTC}


\textbf{Fuente}: configuración de \texttt{@deepseek-ai/dsh-tools}


\begin{lstlisting}
plugins:
  '@deepseek-ai/dsh-tools':
    mode: ptc
    maxParallelSubCalls: 5
\end{lstlisting}


\subsection*{Ejemplo 13: Configuración de Python runtime}


\textbf{Fuente}: configuración de \texttt{@deepseek-ai/dsh-experimental-code-runtime-python}


\begin{lstlisting}
# cordis.yml
'@deepseek-ai/dsh-experimental-code-runtime-python':
  cpuSeconds: 10
  maxWallMs: 30000
  addressSpaceMb: 512
  maxLogBytes: 1048576
  maxValueBytes: 524288
  pythonBin: /usr/bin/python3
\end{lstlisting}


\subsection*{Ejemplo 14: Configuración de headless}


\textbf{Fuente}: configuración de \texttt{@deepseek-ai/dsh-headless}


\begin{lstlisting}
'@deepseek-ai/dsh-headless':
  task: "Analiza el archivo README.md y sugiere mejoras"
\end{lstlisting}


\subsection*{Ejemplo 15: package.json para Remote}


\textbf{Fuente}: patrón de package.json para servicios remotos


\begin{lstlisting}
{
  "exports": {
    "./typert": { "types": "./lib/typert.host.d.ts", "default": "./lib/typert.host.js" },
    "./remote": { "types": "./lib/typert.remote-client.d.ts", "default": "./lib/typert.remote-client.js" }
  },
  "peerDependencies": { "@deepseek-ai/dsh-typert-protocol": "workspace:^" },
  "devDependencies": { "@deepseek-ai/dsh-typert-protocol": "workspace:^" }
}
\end{lstlisting}


\subsection*{Ejemplo 16: Hook de permission gate}


\textbf{Fuente}: patrón de hooks en el repo


\begin{lstlisting}
export function apply(ctx: Context) {
  ctx.on('tools/pre-execute', async (exec, next): Promise<PreToolDecision> => {
    if (!(await isAllowed(exec))) {
      return { kind: 'deny', reason: 'Denied by policy.' }
    }
    return next()
  })
}
\end{lstlisting}


\subsection*{Ejemplo 17: UI — streaming en tiempo real}


\textbf{Fuente}: patrón de UI en el repo


\begin{lstlisting}
export function apply(ctx: Context) {
  ctx.on('agent/assistant-stream', ({ frame }) => {
    if (frame.type === 'chunk' && frame.chunk.type === 'text-delta') {
      render(frame.chunk.text)
    }
  })

  onUserInput(text => ctx.agents.get(sessionId)?.followup(createUserMessage({
    content: [{ type: 'text', text }],
    source: { kind: 'user' },
  })))
}
\end{lstlisting}


\subsection*{Ejemplo 18: Protocol driver (stdio)}


\textbf{Fuente}: patrón de protocol driver en el repo


\begin{lstlisting}
export function apply(ctx: Context) {
  ctx.on('session/event', (_session, event) => {
    if (event.type === 'assistant/message' || event.type === 'assistant/attempt') {
      for (const { chunk } of expandAssistantStream(event.data.stream)) {
        if (chunk.type === 'text-delta') { /* enviar a cliente */ }
      }
    }
  })
  // Entrada: crear/recuperar agente, alimentar con followup(), devolver recibo de encolado
  // Desmontaje: AgentHandle.dispose() (stop + await exit)
}
\end{lstlisting}


\subsection*{Ejemplo 19: Waterfall para decisiones}


\textbf{Fuente}: patrón de waterfall en el repo


\begin{lstlisting}
// Listener que envuelve:
ctx.on('internal/update', function (config, next) {
  console.log('antes:', config)
  const result = next(config)   // delega al siguiente listener / built-in
  console.log('después:', result)
  return result                 // propaga (posiblemente modificado)
})

// Política que veta:
ctx.on('some-decision-event', function (payload, next) {
  if (payload.blocked) {
    return { denied: true }     // cortocircuito – no llama a next()
  }
  return next()
})
\end{lstlisting}


\subsection*{Ejemplo 20: Configuración validada con Schemastery}


\textbf{Fuente}: patrón de configuración en el repo


\begin{lstlisting}
import Schema from '@deepseek-ai/schemastery'

export interface Config {
  greeting: string
  targets: string[]
}

export const Config: Schema<Config> = Schema.object({
  greeting: Schema.string().default('Hello'),
  targets: Schema.array(String).default(['world']),
})

export function apply(ctx: Context, config: Config) {
  for (const target of config.targets) {
    console.log(`${config.greeting}, ${target}!`)
  }
}
\end{lstlisting}


\begin{lstlisting}
# cordis.yml
- name: './config-demo.ts'
  config:
    targets: ['alpha', 'beta']
\end{lstlisting}


\section*{glosario}

\section*{Glosario técnico del proyecto dsh}


\subsection*{A}


\begin{itemize}[leftmargin=*,topsep=2pt,itemsep=1pt]
  \item \textbf{Append-only}: Propiedad de los logs de sesión: son inmutables y solo crecen. Los eventos durables (\texttt{turn/\textit{}, \texttt{step/}}, \texttt{user/message}, \texttt{assistant/message}, \texttt{assistant/attempt}, \texttt{tool/*}) no se modifican ni eliminan.
  \item \textbf{Agent Client Protocol (ACP)}: Protocolo de comunicación con agentes por stdio.
  \item \textbf{AgentHandle}: Handle vivo de un agente. Métodos: \texttt{followup()}, \texttt{steer()}, \texttt{inject()}, \texttt{dispose()}.
\end{itemize}


\subsection*{B}


\begin{itemize}[leftmargin=*,topsep=2pt,itemsep=1pt]
  \item \textbf{Bundle}: Paquete que compone una aplicación concreta (ej. \texttt{ctx.agentLoop}). Los bundles internos son \texttt{dsh-base}, \texttt{dsh-web-app}, \texttt{dsh-headless}, \texttt{dsh-sdk-app}, \texttt{dsh-sdk-minimal}, \texttt{dsh-acp-app}.
  \item \textbf{Bail}: Modo de dispatch de eventos: secuencial síncrono, se detiene en el primer listener que retorna un valor.
\end{itemize}


\subsection*{C}


\begin{itemize}[leftmargin=*,topsep=2pt,itemsep=1pt]
  \item \textbf{Cordis}: Kernel de plugins del harness. Proporciona \texttt{Context}, \texttt{Service}, \texttt{Fiber}, eventos y ciclo de vida.
  \item \textbf{Core}: Servicio único, dueño del estado (ej. \texttt{ctx.sessions}, \texttt{ctx.tools}).
  \item \textbf{ctx}: Objeto central de Cordis: servicios, eventos y ciclo de vida accesibles a través de \texttt{ctx}.
\end{itemize}


\subsection*{D}


\begin{itemize}[leftmargin=*,topsep=2pt,itemsep=1pt]
  \item \textbf{Driver}: Orquestador del bucle del agente. Opera en 4 modos de ejecución en tiempo real.
  \item \textbf{Disposer}: Función cleanup retornada por un efecto. Se ejecuta en orden inverso de registro.
  \item \textbf{DispatchMode}: Tipo de dispatch de eventos: \texttt{'emit' | 'parallel' | 'serial' | 'bail' | 'waterfall'}.
\end{itemize}


\subsection*{E}


\begin{itemize}[leftmargin=*,topsep=2pt,itemsep=1pt]
  \item \textbf{Emit}: Modo de dispatch de eventos: broadcast sin retorno, sin orden garantizado.
  \item \textbf{Effect}: Recurso externo registrado con \texttt{ctx.effect()}. El disposer se ejecuta al descargar el fiber.
  \item \textbf{EffectMeta}: Metadatos de diagnóstico de efectos: \texttt{{ label: string, children: EffectMeta[] }}.
\end{itemize}


\subsection*{F}


\begin{itemize}[leftmargin=*,topsep=2pt,itemsep=1pt]
  \item \textbf{Fiber}: Instancia de plugin en runtime. Estados: \texttt{PENDING → LOADING → ACTIVE → UNLOADING → DISPOSED} (o \texttt{FAILED}).
  \item \textbf{FiberState}: Enum de estados del fiber.
\end{itemize}


\subsection*{G}


\begin{itemize}[leftmargin=*,topsep=2pt,itemsep=1pt]
  \item \textbf{Gateway}: Resuelve objetos complejos del Host (ej. \texttt{Agent}) usando \texttt{ctx.typert.lookups}.
\end{itemize}


\subsection*{H}


\begin{itemize}[leftmargin=*,topsep=2pt,itemsep=1pt]
  \item \textbf{HMR (Hot Module Reload)}: Recarga de plugins al guardar archivos. Requiere \texttt{cordis-plugin-hmr}.
  \item \textbf{Headless}: Modo de ejecución one-shot sin interfaz. \texttt{dsh --profile headless "task"}.
\end{itemize}


\subsection*{I}


\begin{itemize}[leftmargin=*,topsep=2pt,itemsep=1pt]
  \item \textbf{Inject}: Declaración de dependencias duras de un plugin. \texttt{inject: ['servicio1', 'servicio2']}.
  \item \textbf{Isolate}: Crear un hijo con ámbito de servicio independiente para un nombre.
\end{itemize}


\subsection*{J}


\begin{itemize}[leftmargin=*,topsep=2pt,itemsep=1pt]
  \item \textbf{JSON-RPC}: Protocolo RPC por stdio usado por los perfiles \texttt{sdk} y \texttt{sdk-minimal}.
\end{itemize}


\subsection*{L}


\begin{itemize}[leftmargin=*,topsep=2pt,itemsep=1pt]
  \item \textbf{Loader}: Cargador de plugins. Resuelve bundles y plugins desde la instalación de \texttt{dsh} o \texttt{node\_modules} del perfil.
  \item \textbf{Live reload}: Modo de recarga de parches que observa cambios en archivos y reaplica (\texttt{patchReload: "live"}).
\end{itemize}


\subsection*{M}


\begin{itemize}[leftmargin=*,topsep=2pt,itemsep=1pt]
  \item \textbf{MCP (Model Context Protocol)}: Protocolo para conectar servidores de herramientas. Transportes: \texttt{stdio} y \texttt{streamable-http}.
  \item \textbf{Monotonic guard}: Guardia monotónica para denegación final en el pipeline de herramientas.
\end{itemize}


\subsection*{O}


\begin{itemize}[leftmargin=*,topsep=2pt,itemsep=1pt]
  \item \textbf{Overlay}: Capa de parche aplicada con \texttt{--patch <path>}. Las capas posteriores sobrescriben filas completas.
\end{itemize}


\subsection*{P}


\begin{itemize}[leftmargin=*,topsep=2pt,itemsep=1pt]
  \item \textbf{Parallel}: Modo de dispatch de eventos: ejecución concurrente de todos los listeners.
  \item \textbf{Patch}: Archivo YAML que modifica la configuración. Capas: bundles → perfil → home → overlays.
  \item \textbf{Perfil}: Composición de bundles y parches. Se inicia con \texttt{dsh --profile <nombre>}.
  \item \textbf{PTC (Prompt-to-Code)}: Modo de presentación de herramientas donde el código se genera como herramienta.
  \item \textbf{Plugin}: Unidad de extensión Cordis. Formas: función, clase con \texttt{apply}, u objeto \texttt{{ apply }}.
\end{itemize}


\subsection*{R}


\begin{itemize}[leftmargin=*,topsep=2pt,itemsep=1pt]
  \item \textbf{Remote}: Decorador para exponer un método como endpoint RPC unario.
  \item \textbf{RemoteScope}: Decorador para métodos que operan sobre un contexto resuelto por ámbito (ej. \texttt{'agent'}).
  \item \textbf{RemoteError}: Error tipado para el cliente remoto. Códigos declarados en \texttt{RemoteErrorDetailsMap}.
\end{itemize}


\subsection*{S}


\begin{itemize}[leftmargin=*,topsep=2pt,itemsep=1pt]
  \item \textbf{Seam}: Interfaz abstracta (ej. \texttt{ctx.llm}, \texttt{ctx.fs}). Varios backends pueden registrarse.
  \item \textbf{Serial}: Modo de dispatch de eventos: ejecución uno a uno, espera cada listener.
  \item \textbf{Service}: Clase base para servicios en \texttt{ctx.<name>}. Registra la instancia en el contexto.
  \item \textbf{Schemastery}: Librería de validación de esquemas de configuración.
  \item \textbf{Settlement}: Frame de fin confirmado que consolida un stream como \texttt{assistant/message} o \texttt{assistant/attempt}.
  \item \textbf{Spill}: Desbordamiento de resultados de herramientas a archivos cuando exceden \texttt{maxInlineBytes}.
  \item \textbf{Subagent}: Agente hijo delegado. Proveedores: \texttt{spawn}, \texttt{fork}, \texttt{acp}, \texttt{codex}, \texttt{claude-code}, \texttt{dsh-sdk}.
\end{itemize}


\subsection*{T}


\begin{itemize}[leftmargin=*,topsep=2pt,itemsep=1pt]
  \item \textbf{Typert}: Protocolo RPC tipado entre Host y Cliente. Artefactos: \texttt{typert.host.js}, \texttt{typert.remote-client.js}.
  \item \textbf{TypertRemoteService}: Clase base para servicios de negocio expuestos al cliente.
  \item \textbf{Tool}: Herramienta registrada en \texttt{ctx.tools}. Pipeline: pre-policy → monotonic guards → dispatch → post-policy → final-result.
  \item \textbf{ToolCallId}: Identificador de llamada a herramienta con marca de tipo.
\end{itemize}


\subsection*{W}


\begin{itemize}[leftmargin=*,topsep=2pt,itemsep=1pt]
  \item \textbf{Waterfall}: Modo de dispatch de eventos: middleware encadenado. Cada listener recibe \texttt{(...args, next)}.
  \item \textbf{Workspace}: Entidad de espacio de trabajo con ID de marca, registrada en \texttt{ctx.workspaceRegistry}.
\end{itemize}


\subsection*{Y}


\begin{itemize}[leftmargin=*,topsep=2pt,itemsep=1pt]
  \item \textbf{YAML \texttt{!!js}}: Expresión JavaScript evaluada en tiempo de carga dentro de \texttt{config} y \texttt{disabled}.
\end{itemize}


\section*{patrones}

\section*{Patrones de uso comunes}


\subsection*{1. Envolver recurso externo en un efecto}


\begin{lstlisting}
export function apply(ctx: Context) {
  ctx.effect(() => {
    const timer = setInterval(() => console.log('tick'), 200)
    return () => clearInterval(timer)   // disposer
  })
}
\end{lstlisting}


\textbf{Flujo}:


\begin{enumerate}[leftmargin=*,topsep=2pt,itemsep=1pt]
  \item \texttt{apply} se ejecuta al cargar el plugin.
  \item \texttt{ctx.effect} ejecuta el callback inmediatamente.
  \item El disposer retornado se registra para cleanup.
  \item Al descargar el fiber, el disposer se ejecuta en orden inverso.
\end{enumerate}


\subsection*{2. Montar un sub-plugin y desmontarlo luego}


\begin{lstlisting}
const fiber = ctx.plugin(heartbeat)
// ... más tarde
await fiber.dispose()
\end{lstlisting}


\textbf{Flujo}:


\begin{enumerate}[leftmargin=*,topsep=2pt,itemsep=1pt]
  \item \texttt{ctx.plugin} carga el plugin y retorna un Fiber awaitable.
  \item El fiber pasa por estados: PENDING → LOADING → ACTIVE.
  \item \texttt{dispose()} ejecuta cleanup y pasa a DISPOSED.
\end{enumerate}


\subsection*{3. Proveer un servicio}


\begin{lstlisting}
declare module '@deepseek-ai/cordis' {
  interface Context { greeter: GreeterService }
}

export class GreeterService extends Service {
  constructor(ctx: Context) { super(ctx, 'greeter') }
  greet(who: string) { return `Hello, ${who}!` }
}

export function apply(ctx: Context) { ctx.plugin(GreeterService) }
\end{lstlisting}


\textbf{Flujo}:


\begin{enumerate}[leftmargin=*,topsep=2pt,itemsep=1pt]
  \item Declarar el servicio en la interfaz \texttt{Context} (module augmentation).
  \item Extender \texttt{Service} y llamar \texttt{super(ctx, 'greeter')} en el constructor.
  \item Registrar con \texttt{ctx.plugin(GreeterService)}.
  \item Otros plugins pueden consumir con \texttt{inject: ['greeter']}.
\end{enumerate}


\subsection*{4. Consumir un servicio con dependencia dura}


\begin{lstlisting}
export const inject = ['greeter']
export function apply(ctx: Context) {
  console.log(ctx.greeter.greet('world'))
}
\end{lstlisting}


\textbf{Flujo}:


\begin{enumerate}[leftmargin=*,topsep=2pt,itemsep=1pt]
  \item Declarar \texttt{inject} con los nombres de servicios requeridos.
  \item El plugin espera a que los servicios estén disponibles.
  \item Si un servicio requerido cambia, el plugin se recarga automáticamente.
\end{enumerate}


\subsection*{5. Consumir un servicio opcional}


\begin{lstlisting}
export function apply(ctx: Context) {
  const greeter = ctx.get('greeter')
  console.log(greeter?.greet('maybe') ?? 'no greeter available')
}
\end{lstlisting}


\textbf{Flujo}:


\begin{enumerate}[leftmargin=*,topsep=2pt,itemsep=1pt]
  \item Usar \texttt{ctx.get('greeter')} sin declarar \texttt{inject}.
  \item Si el servicio no existe, retorna \texttt{undefined}.
  \item Manejar el caso con optional chaining y fallback.
\end{enumerate}


\subsection*{6. Declarar y emitir eventos tipados}


\begin{lstlisting}
declare module '@deepseek-ai/cordis' {
  interface Events {
    'stats/report'(name: string, count: number): void
  }
}

// Emisión
this.ctx.emit('stats/report', name, next)

// Escucha (efecto automático)
ctx.on('stats/report', (name, count) => {
  console.log(`[stats] ${name} -> ${count}`)
})
\end{lstlisting}


\textbf{Flujo}:


\begin{enumerate}[leftmargin=*,topsep=2pt,itemsep=1pt]
  \item Declarar el evento en la interfaz \texttt{Events} (module augmentation).
  \item Emitir con \texttt{ctx.emit}.
  \item Escuchar con \texttt{ctx.on} — el listener se registra como efecto y se limpia automáticamente.
\end{enumerate}


\subsection*{7. Waterfall: transformar o cortocircuitar}


\begin{lstlisting}
declare module '@deepseek-ai/cordis' {
  interface Events {
    'demo/transform'(input: string, next: () => Promise<string>): Promise<string>
  }
}

// Listener que transforma
ctx.on('demo/transform', async (input, next) => {
  const downstream = await next()
  return downstream.toUpperCase()
})

// Listener que cortocircuita
ctx.on('demo/transform', async (input, next) => {
  if (input.includes('blocked')) return '** blocked **'
  return next()
})

// Invocación
const result = await ctx.waterfall('demo/transform', 'hello', async () => 'hello')
\end{lstlisting}


\textbf{Flujo}:


\begin{enumerate}[leftmargin=*,topsep=2pt,itemsep=1pt]
  \item Cada listener recibe \texttt{(input, next)}.
  \item \texttt{next()} invoca el siguiente listener en la cadena.
  \item No llamar a \texttt{next()} cortocircuita la cadena.
  \item El valor de retorno se propaga hacia afuera.
\end{enumerate}


\subsection*{8. Configuración validada con Schemastery}


\begin{lstlisting}
import Schema from '@deepseek-ai/schemastery'

export interface Config {
  greeting: string
  targets: string[]
}

export const Config: Schema<Config> = Schema.object({
  greeting: Schema.string().default('Hello'),
  targets: Schema.array(String).default(['world']),
})

export function apply(ctx: Context, config: Config) {
  for (const target of config.targets) {
    console.log(`${config.greeting}, ${target}!`)
  }
}
\end{lstlisting}


\begin{lstlisting}
# cordis.yml
- name: './config-demo.ts'
  config:
    targets: ['alpha', 'beta']
\end{lstlisting}


\textbf{Flujo}:


\begin{enumerate}[leftmargin=*,topsep=2pt,itemsep=1pt]
  \item Definir interfaz \texttt{Config} y schema Schemastery.
  \item El loader valida la config del YAML contra el schema.
  \item \texttt{apply} recibe la config validada.
\end{enumerate}


\subsection*{9. Valores computados en YAML (\texttt{!!js})}


\begin{lstlisting}
- name: './config-demo.ts'
  config:
    greeting: !!js process.env.DEMO_GREETING ?? 'Hello'
\end{lstlisting}


\textbf{Flujo}:


\begin{enumerate}[leftmargin=*,topsep=2pt,itemsep=1pt]
  \item \texttt{!!js} evalúa la expresión JavaScript en tiempo de carga.
  \item Solo permitido dentro de \texttt{config} y \texttt{disabled}.
  \item En modo \texttt{live}, editar el archivo reevalúa expresiones.
\end{enumerate}


\subsection*{10. Diagnóstico de plugins en PENDING}


\begin{lstlisting}
import { FiberState, type Context } from '@deepseek-ai/cordis'

export function apply(ctx: Context) {
  setTimeout(() => {
    for (const runtime of ctx.registry.values()) {
      for (const fiber of runtime.fibers) {
        if (fiber.state === FiberState.PENDING) {
          console.log(`${fiber.name} is PENDING — a required service is missing`)
        }
      }
    }
  }, 500)
}
\end{lstlisting}


\textbf{Flujo}:


\begin{enumerate}[leftmargin=*,topsep=2pt,itemsep=1pt]
  \item Iterar \texttt{ctx.registry.values()} para obtener runtimes.
  \item Iterar \texttt{runtime.fibers} para obtener fibers.
  \item Verificar \texttt{fiber.state === FiberState.PENDING}.
\end{enumerate}


\subsection*{11. Servicio remoto Typert (Host)}


\begin{lstlisting}
import type { Agent } from '@deepseek-ai/dsh-agent'
import { TypertRemoteService, Remote, RemoteScope } from '@deepseek-ai/dsh-typert-protocol'
import type { Context } from '@deepseek-ai/cordis'

export class GoalService extends TypertRemoteService {
  constructor(ctx: Context) {
    super(ctx, 'goals')
  }

  @Remote('create')
  createForClient(
    agent: Agent,
    request: { objective: string },
    signal: AbortSignal,
  ): { accepted: boolean } {
    signal.throwIfAborted()
    return this.create(agent, request)
  }

  @RemoteScope('agent', 'current')
  currentForClient(): { accepted: boolean } {
    return { accepted: true }
  }

  private create(_agent: Agent, request: { objective: string }): { accepted: boolean } {
    return { accepted: request.objective.length > 0 }
  }
}
\end{lstlisting}


\textbf{Flujo}:


\begin{enumerate}[leftmargin=*,topsep=2pt,itemsep=1pt]
  \item Extender \texttt{TypertRemoteService} con namespace \texttt{'goals'}.
  \item Decorar métodos públicos con \texttt{@Remote} o \texttt{@RemoteScope}.
  \item El Gateway resuelve objetos complejos (ej. \texttt{Agent}) usando \texttt{ctx.typert.lookups}.
\end{enumerate}


\subsection*{12. Consumo remoto con manejo de errores (Cliente)}


\begin{lstlisting}
import type { Context } from '@deepseek-ai/cordis'
import { isRemoteFailure } from '@deepseek-ai/dsh-api-gateway/client'
import type {} from '@deepseek-ai/dsh-api-remotes/client'

export const inject = ['remote', 'remote.notes']

declare const ctx: Context

export async function noteTitles(): Promise<readonly string[]> {
  const result = await ctx.remote.notes.list()
  if (!result.ok) {
    if (result.error.code === 'note/not-found') return []
    throw result.error
  }
  return result.value.map(row => row.title)
}

export async function renderTitles(): Promise<string> {
  try {
    return (await noteTitles()).join(', ')
  } catch (error: unknown) {
    if (!isRemoteFailure(error)) throw error
    return `unavailable (${error.code})`
  }
}

export function hostLabel(): string {
  const { home, isLoopback } = ctx.remote.$host
  return home ?? (isLoopback ? 'local host' : 'remote host')
}
\end{lstlisting}


\textbf{Flujo}:


\begin{enumerate}[leftmargin=*,topsep=2pt,itemsep=1pt]
  \item Declarar \texttt{inject: ['remote', 'remote.notes']}.
  \item Llamar \texttt{ctx.remote.notes.list()}.
  \item Verificar \texttt{result.ok}; si es \texttt{false}, ramificar por \texttt{result.error.code}.
  \item Usar \texttt{isRemoteFailure(error)} para distinguir fallos remotos de defectos locales.
\end{enumerate}


\subsection*{13. Hook de permission gate}


\begin{lstlisting}
export function apply(ctx: Context) {
  ctx.on('tools/pre-execute', async (exec, next): Promise<PreToolDecision> => {
    if (!(await isAllowed(exec))) {
      return { kind: 'deny', reason: 'Denied by policy.' }
    }
    return next()
  })
}
\end{lstlisting}


\textbf{Flujo}:


\begin{enumerate}[leftmargin=*,topsep=2pt,itemsep=1pt]
  \item Escuchar \texttt{tools/pre-execute} (waterfall).
  \item Si la política deniega, retornar \texttt{{ kind: 'deny', reason }} sin llamar a \texttt{next()}.
  \item Si permite, llamar a \texttt{next()} para continuar la cadena.
\end{enumerate}


\subsection*{14. UI: streaming en tiempo real}


\begin{lstlisting}
export function apply(ctx: Context) {
  ctx.on('agent/assistant-stream', ({ frame }) => {
    if (frame.type === 'chunk' && frame.chunk.type === 'text-delta') {
      render(frame.chunk.text)
    }
  })

  onUserInput(text => ctx.agents.get(sessionId)?.followup(createUserMessage({
    content: [{ type: 'text', text }],
    source: { kind: 'user' },
  })))
}
\end{lstlisting}


\textbf{Flujo}:


\begin{enumerate}[leftmargin=*,topsep=2pt,itemsep=1pt]
  \item Escuchar \texttt{agent/assistant-stream} para chunks transitorios.
  \item Enviar entrada de usuario con \texttt{ctx.agents.get(sessionId)?.followup(...)}.
\end{enumerate}


\subsection*{15. Protocol driver (stdio)}


\begin{lstlisting}
export function apply(ctx: Context) {
  ctx.on('session/event', (_session, event) => {
    if (event.type === 'assistant/message' || event.type === 'assistant/attempt') {
      for (const { chunk } of expandAssistantStream(event.data.stream)) {
        if (chunk.type === 'text-delta') { /* enviar a cliente */ }
      }
    }
  })
  // Entrada: crear/recuperar agente, alimentar con followup(), devolver recibo de encolado
  // Desmontaje: AgentHandle.dispose() (stop + await exit)
}
\end{lstlisting}


\textbf{Flujo}:


\begin{enumerate}[leftmargin=*,topsep=2pt,itemsep=1pt]
  \item Escuchar \texttt{session/event} para eventos durables.
  \item Expandir el stream del asistente y enviar deltas de texto al cliente.
  \item Para entrada, crear/recuperar agente y usar \texttt{followup()}.
\end{enumerate}


\subsection*{16. Pipeline de herramientas}


\begin{lstlisting}
// ctx.tools expone:
//   - Registro de capacidades
//   - Transporte PTC mode
//   - Ruteo con pre-policy → monotonic guards → dispatch → post-policy → final-result

// Herramientas registradas:
ctx.tools.register(toolBash)     // usa ctx.shell, ctx.approval, ctx.jobs
ctx.tools.register(toolFs)       // usa ctx.fs
ctx.tools.register(toolWeb)      // usa ctx.web
ctx.tools.register(toolSubagent) // usa ctx.subagents
\end{lstlisting}


\textbf{Flujo}:


\begin{enumerate}[leftmargin=*,topsep=2pt,itemsep=1pt]
  \item Registrar herramientas con \texttt{ctx.tools.register}.
  \item El pipeline ejecuta: pre-policy → monotonic guards → dispatch → post-policy → final-result.
  \item Eventos: \texttt{tools/pre-execute} → \texttt{tools/execute} → \texttt{tools/post-execute} → \texttt{tools/result}.
\end{enumerate}


\subsection*{17. Flujo del turno del agente}


\begin{lstlisting}
turn/start → (pasos) → agent/turn-stopping → turn/end
\end{lstlisting}


\textbf{Modos de ejecución}:


\begin{enumerate}[leftmargin=*,topsep=2pt,itemsep=1pt]
  \item \textbf{Inbox/Cola}: El \texttt{Driver} se activa con trabajo en cola. Reclama mensajes de entrada (\texttt{agent/inbox/claimed}).
  \item \textbf{Pre-paso (Waterfall)}: Emite \texttt{agent/pre-step}. Los hooks pueden rechazar autoritativamente o ingresar (\texttt{enter(messages)}) el lote propuesto. Un rechazo cierra el turno sin ejecutar paso.
  \item \textbf{Paso de ejecución}: Si se ingresa, inicia \texttt{step/start}. Ensambla el prompt del sistema (\texttt{system-prompt/assemble}), realiza la solicitud al LLM (\texttt{agent/request} → \texttt{llm/stream}), y transmite los chunks de vuelta (\texttt{agent/assistant-stream}). El resultado exitoso se registra como \texttt{assistant/message}; los fallos/errores como \texttt{assistant/attempt}.
  \item \textbf{Ejecución de herramientas}: Clasifica y ejecuta llamadas a herramientas en un pool rodante con barreras. Emite \texttt{tool/call} y \texttt{tool/result}. Si las herramientas generan más trabajo, se inicia un nuevo paso.
\end{enumerate}


\subsection*{18. Trazabilidad append-only}


\begin{itemize}[leftmargin=*,topsep=2pt,itemsep=1pt]
  \item \textbf{Eventos de sesión}: \texttt{turn/\textit{}, \texttt{step/}}, \texttt{user/message}, \texttt{assistant/message}, \texttt{assistant/attempt}, y \texttt{tool/*} son eventos durables e inmutables en el log de sesión.
  \item \textbf{Transmisión en tiempo real}: \texttt{agent/assistant-stream} publica chunks transitorios. El loop consolida el stream completo como un \texttt{assistant/message} o \texttt{assistant/attempt} antes de un frame de fin confirmado.
  \item \textbf{Recuperabilidad}: Todo lo que el modelo ve debe ser reconstruible desde el log. Si un proceso se interrumpe antes de un \texttt{settlement}, no queda un stream de intento persistente. Los consumidores de transcripciones deben usar \texttt{session/event}.
\end{itemize}


\subsection*{19. Almacenamiento con backends múltiples}


\begin{lstlisting}
// Backends nombrados bajo ctx.storage
storage.register('json', storageJson)
storage.register('sqlite', storageSqlite)

// ctx.storageDomain espera a cada backend configurado y publica
// un único servicio tipado para estado duradero
\end{lstlisting}


\textbf{Flujo}:


\begin{enumerate}[leftmargin=*,topsep=2pt,itemsep=1pt]
  \item Registrar backends con \texttt{storage.register(nombre, implementación)}.
  \item \texttt{ctx.storageDomain} enruta dominios a backends específicos.
  \item Configuración: \texttt{backend} (default) y \texttt{routes} (overrides por dominio).
\end{enumerate}


\subsection*{20. Sesiones y persistencia}


\begin{lstlisting}
// ctx.sessions: append-only, in-memory
// Emite feed de eventos durables

// ctx.sessionPersistence (seam) → backends reales
// Implementación JSONL: un artifact por Session con SessionEvent vocabulary
// Consumidores: agent-loop, tool-bash, hooks, session-query, message-feedback

// ctx.sessionProjectionCache: checkpoint cache + cold-read ladder
// Puntos mandatorios: throttled + turn/end/detach
\end{lstlisting}


\textbf{Flujo}:


\begin{enumerate}[leftmargin=*,topsep=2pt,itemsep=1pt]
  \item \texttt{ctx.sessions} mantiene el log en memoria.
  \item \texttt{ctx.sessionPersistence} persiste a JSONL.
  \item \texttt{ctx.sessionProjectionCache} mantiene checkpoints para cold-read.
\end{enumerate}


\section*{prerrequisitos}

\section*{Prerrequisitos del proyecto dsh}


\subsection*{Dependencias del sistema}


\begin{itemize}[leftmargin=*,topsep=2pt,itemsep=1pt]
  \item \textbf{Node.js}: runtime JavaScript/TypeScript.
  \item \textbf{pnpm}: gestor de paquetes. Debe estar en el PATH para \texttt{dsh plugin}.
  \item \textbf{pnpm ≥10}: para plugins desde Git con \texttt{prepare} script, bloquea hasta permitir en \texttt{pnpm-workspace.yaml}.
  \item \textbf{CPython 3.10+}: requerido por \texttt{@deepseek-ai/dsh-experimental-code-runtime-python}.
\end{itemize}


\subsection*{Paquetes npm principales}


\subsubsection*{Core}


\begin{itemize}[leftmargin=*,topsep=2pt,itemsep=1pt]
  \item \texttt{@deepseek-ai/cordis} — kernel de plugins (Context, Service, Fiber, eventos).
  \item \texttt{@deepseek-ai/dsh-agent} — fábrica de agentes.
  \item \texttt{@deepseek-ai/dsh-agent-loop} — driver del bucle agente.
  \item \texttt{@deepseek-ai/dsh-session} — almacén de sesiones append-only.
  \item \texttt{@deepseek-ai/dsh-tools} — registro de herramientas y pipeline.
  \item \texttt{@deepseek-ai/dsh-system-prompt} — ensamblador de system prompt.
\end{itemize}


\subsubsection*{Protocolo}


\begin{itemize}[leftmargin=*,topsep=2pt,itemsep=1pt]
  \item \texttt{@deepseek-ai/dsh-typert-protocol} — RPC tipado entre Host y Cliente.
  \item \texttt{@deepseek-ai/dsh-api-gateway} — gateway de API.
  \item \texttt{@deepseek-ai/dsh-api-remotes} — remotes de API.
\end{itemize}


\subsubsection*{Bundles}


\begin{itemize}[leftmargin=*,topsep=2pt,itemsep=1pt]
  \item \texttt{@deepseek-ai/dsh-base}
  \item \texttt{@deepseek-ai/dsh-web-app}
  \item \texttt{@deepseek-ai/dsh-headless}
  \item \texttt{@deepseek-ai/dsh-sdk-app}
  \item \texttt{@deepseek-ai/dsh-sdk-minimal}
  \item \texttt{@deepseek-ai/dsh-acp-app}
\end{itemize}


\subsubsection*{Seams y backends}


\begin{itemize}[leftmargin=*,topsep=2pt,itemsep=1pt]
  \item \texttt{@deepseek-ai/dsh-llm} + \texttt{dsh-llm-deepseek}, \texttt{dsh-llm-pi-ai}, \texttt{dsh-llm-replay}
  \item \texttt{@deepseek-ai/dsh-fs} + \texttt{dsh-fs-local}, \texttt{dsh-fs-sandbox}, \texttt{dsh-fs-e2b}
  \item \texttt{@deepseek-ai/dsh-storage} + \texttt{dsh-storage-json}, \texttt{dsh-storage-sqlite}
  \item \texttt{@deepseek-ai/dsh-subagents} + \texttt{dsh-subagent-spawn-in-process}, \texttt{dsh-subagent-fork-in-process}, \texttt{dsh-subagent-acp}, \texttt{dsh-subagent-codex}, \texttt{dsh-subagent-claude-code}, \texttt{dsh-subagent-dsh-sdk}
  \item \texttt{@deepseek-ai/dsh-compaction} + \texttt{dsh-compaction-basic}
  \item \texttt{@deepseek-ai/dsh-sandbox} + \texttt{dsh-sandbox-local}
  \item \texttt{@deepseek-ai/dsh-shell} + \texttt{dsh-bash-local}, \texttt{dsh-bash-sandbox}, \texttt{dsh-pwsh-local}
  \item \texttt{@deepseek-ai/dsh-web} + \texttt{dsh-web-search-exa}, \texttt{dsh-web-search-perplexity}, \texttt{dsh-web-search-deepseek}, \texttt{dsh-web-fetch-http}
\end{itemize}


\subsubsection*{Herramientas}


\begin{itemize}[leftmargin=*,topsep=2pt,itemsep=1pt]
  \item \texttt{@deepseek-ai/dsh-tool-bash}
  \item \texttt{@deepseek-ai/dsh-tool-fs}
  \item \texttt{@deepseek-ai/dsh-tool-web}
  \item \texttt{@deepseek-ai/dsh-tool-subagent}
  \item \texttt{@deepseek-ai/dsh-tool-terminal}
  \item \texttt{@deepseek-ai/dsh-tool-jobs}
  \item \texttt{@deepseek-ai/dsh-tool-workflow}
  \item \texttt{@deepseek-ai/dsh-tool-lsp}
  \item \texttt{@deepseek-ai/dsh-tool-skill}
  \item \texttt{@deepseek-ai/dsh-tool-cordis}
  \item \texttt{@deepseek-ai/dsh-tool-ralph}
  \item \texttt{@deepseek-ai/dsh-tool-todo}
\end{itemize}


\subsubsection*{Frontend}


\begin{itemize}[leftmargin=*,topsep=2pt,itemsep=1pt]
  \item \texttt{@deepseek-ai/dsh-host-webserver}
  \item \texttt{@deepseek-ai/dsh-client-connection}
  \item \texttt{@deepseek-ai/dsh-client-modules}
  \item \texttt{@deepseek-ai/dsh-client-hmr}
  \item \texttt{@deepseek-ai/dsh-client-ui-*} (múltiples paquetes UI)
\end{itemize}


\subsubsection*{Testing}


\begin{itemize}[leftmargin=*,topsep=2pt,itemsep=1pt]
  \item \texttt{@deepseek-ai/dsh-client-test-runtime} — \texttt{TestRemote}, \texttt{RemoteError} para pruebas de cliente.
  \item \texttt{vitest} — framework de pruebas.
\end{itemize}


\subsubsection*{Utilidades}


\begin{itemize}[leftmargin=*,topsep=2pt,itemsep=1pt]
  \item \texttt{@deepseek-ai/schemastery} — validación de esquemas de configuración.
  \item \texttt{@deepseek-ai/dsh-brand} — utilidades de branding (ej. \texttt{brandString}).
\end{itemize}


\subsection*{Instalación}


\begin{lstlisting}
# Instalar dsh globalmente (o usar pnpm dlx)
pnpm add -g @deepseek-ai/dsh

# Verificar instalación
dsh --version

# Inicializar un perfil
dsh --profile web

# Añadir plugins a un perfil
dsh plugin --profile <name> add <package>
\end{lstlisting}


\subsection*{Requisitos de configuración}


\begin{itemize}[leftmargin=*,topsep=2pt,itemsep=1pt]
  \item \texttt{$DSH\_HOME}: directorio home del harness (default: \texttt{\textasciitilde{}/.dsh}).
  \item \texttt{$DSH\_AGENTS\_HOME}: directorio de agentes (default: \texttt{\textasciitilde{}/.agents}).
  \item \texttt{pnpm-workspace.yaml}: necesario para plugins desde Git con \texttt{prepare} script.
  \item \texttt{cordis.patch.yml}: archivo de parches del perfil o home.
\end{itemize}


\subsection*{Notas importantes}


\begin{itemize}[leftmargin=*,topsep=2pt,itemsep=1pt]
  \item Todas las configuraciones son validadas al inicio; los errores de configuración provocan fallos ruidosos.
  \item Los tiempos de gracia (\texttt{disposeGraceMs}, \texttt{shutdownTimeoutMs}) están acotados por \texttt{MAX\_TIMER\_DELAY\_MS}.
  \item Las claves de entorno en subagentes se superponen a un entorno padre limpio de credenciales.
  \item El fallback SRC para desarrollo en fuente es más débil y no genera artefactos estrictos de Typert.
  \item Solo los paquetes que llaman a \texttt{ctx.remote.<namespace>} deben declarar la dependencia \texttt{'remote.<namespace>'} en su array \texttt{inject}.
\end{itemize}


\end{document}
